\documentclass[a4paper,12pt,twoside]{article}
\usepackage[top=2.5cm,bottom=2.5cm,inner=3cm,outer=2cm, headheight=1.25cm, footskip=1.25cm]{geometry}
\usepackage{setspace}
\usepackage{fancyhdr}
\usepackage{titlesec}
\usepackage{enumitem}
\usepackage{hyperref}
\usepackage[polish]{babel}
\usepackage{pdfpages} % Dodaj pakiet pdfpages
\usepackage{graphicx}
\usepackage{float}
\usepackage[normalem]{ulem}
\usepackage[nottoc]{tocbibind}

%\hypersetup{colorlinks=true}

\usepackage[T1]{fontenc} % Użyj kodowania czcionek T1
\usepackage[utf8]{inputenc} % Użyj kodowania wejściowego UTF-8 dla znaków specjalnych
\usepackage{helvet} % Czcionka Helvetica
\renewcommand{\familydefault}{\sfdefault} % Ustaw domyślną czcionkę na bezszeryfową
\usepackage{amsmath} % Potrzebne do korzystania z środowisk matematycznych
\usepackage{chngcntr}
\counterwithin{figure}{section} % Numeracja rysunków zgodnie z numerem rozdziału
\counterwithin{table}{section}
\usepackage{amsmath}
\usepackage{tocloft}
\usepackage{nomencl}
\makenomenclature
\newcommand{\listequationsname}{\hspace{-1.0cm}\Large{Spis równań}}
\newlistof{myequations}{equ}{\listequationsname}
\newcommand{\myequations}[1]{%
    \addcontentsline{equ}{myequations}{\protect\numberline{\theequation}#1}\par
}

\usepackage{caption}
\usepackage{xcolor}
\usepackage{booktabs}
\usepackage{array}
\usepackage{tabularx}
\usepackage{tikz}
\usetikzlibrary{positioning,arrows.meta,fit,backgrounds,calc}
\captionsetup[figure]{name=Rys.}

\usepackage{listings}
\lstset{
    language=Python, % Wybierz język programowania
    basicstyle=\ttfamily, % Styl podstawowy
    numbers=left, % Numery linii po lewej stronie
    numberstyle=\tiny\color{gray}, % Styl numerów linii
    frame=single, % Rama wokół kodu
    breaklines=true, % Automatyczne łamanie linii
    captionpos=b, % Pozycja podpisu
    showstringspaces=false, % Nie pokazuj spacji w napisach
    literate={ą}{{\k{a}}}1 {ć}{{\'c}}1 {ę}{{\k{e}}}1 {ł}{{\l{}}}1 {ń}{{\'n}}1 {ó}{{\'o}}1 {ś}{{\'s}}1 {ź}{{\'z}}1 {ż}{{\.z}}1
             {Ą}{{\k{A}}}1 {Ć}{{\'C}}1 {Ę}{{\k{E}}}1 {Ł}{{\L{}}}1 {Ń}{{\'N}}1 {Ó}{{\'O}}1 {Ś}{{\'S}}1 {Ź}{{\'Z}}1 {Ż}{{\.Z}}1,
}

\usepackage{lmodern}
\renewcommand*\familydefault{\sfdefault}

\pagestyle{fancy}
\fancyhead{}
\fancyhead[LO,RE]{}
\fancyfoot{}
\fancyfoot[LE,RO]{\thepage}
\fancyfoot[LO,RE]{}
\renewcommand{\headrulewidth}{0pt}

\setlength{\parindent}{0.5cm}
\setlist[itemize]{label=•,itemsep=0pt}

\titleformat{\section}{\normalfont\fontsize{16}{18}\bfseries}{\thesection}{1em}{}
\titleformat{\subsection}{\normalfont\fontsize{14}{16}\bfseries}{\thesubsection}{1em}{}
\titleformat{\subsubsection}{\normalfont\fontsize{13}{15}\bfseries}{\thesubsubsection}{1em}{}

\newcommand{\MMp}[1]{\marginpar{\textcolor{blue}{\textbf{MM}: \footnotesize #1}}}
\newcommand{\MM}[1]{\textcolor{blue}{\textbf{MM}: #1}}

\usepackage[backend=bibtex,style=authoryear,natbib]{biblatex}
\addbibresource{literatura.bib}

\fancypagestyle{plain}{
    \fancyhf{} % Usuń wszystko z nagłówka i stopki
    \fancyfoot[LE,RO]{\thepage} % Dodaj numer strony na zewnętrznych krawędziach stopki
    \renewcommand{\headrulewidth}{0pt} % Brak linii w nagłówku
    \renewcommand{\footrulewidth}{0pt} % Brak linii w stopce
}

\fancypagestyle{empty}{
    \fancyhf{} % Usuń wszystko z nagłówka i stopki
    \renewcommand{\headrulewidth}{0pt} % Brak linii w nagłówku
    \renewcommand{\footrulewidth}{0pt} % Brak linii w stopce
}



\begin{document}
\counterwithin{lstlisting}{section}
\thispagestyle{empty}

\includegraphics[width=5in]{LogoPL-FTIMS.pdf}

\begin{centering}

\vspace{1cm}

\textbf{\textit{\large{Igor Ordecha}}}

\vspace{.1cm}

\textbf{\textit{\large{251601}}}

\vspace{1.5cm}

{\large{PRACA DYPLOMOWA}}

\vspace{.1cm}

{\large{inżynierska}}

\vspace{.1cm}

na kierunku Informatyka Stosowana

\vspace{1.8cm}

\textbf{\Large{Weryfikacja skuteczności agenta przeglądarkowego Google Antigravity przy realizacji zadań w~systemach biurowych}}

\end{centering}

\vspace{2.25cm}
\begin{centering}
{Instytut Informatyki I72}\\
\end{centering}
\vspace{.25cm}
\textbf{Promotor:} \hfill dr inż.~Łukasz Chomątek \hfill\mbox{}

\begin{centering}
(tytuł/stopień naukowy, imię i nazwisko)

~\\

\end{centering}
\textbf{Opiekun pomocniczy:*)} \dotfill

\begin{centering}
(tytuł/stopień naukowy, imię i nazwisko)

~\\

\end{centering}
\textbf{Promotor uczelni partnerskiej:**)} \dotfill

\begin{centering}
(tytuł/stopień naukowy, imię i nazwisko)

~\\

\end{centering}

\vfill

\begin{centering}
ŁÓDŹ 2027

\end{centering}
\hspace{0.5cm}
\begin{itemize}
\small
\item[*]	jeśli został powołany
\item[**]	w przypadku procedury uznania
\end{itemize}
 % Wypełnij w tym pliku dane strony tytułowej
%\pagenumbering{roman}
%\setcounter{page}{2}
\newpage
\thispagestyle{empty}
\mbox{}
\newpage
\onehalfspacing

% Spis treści
\pagestyle{empty}
\tableofcontents
\newpage
\clearpage

\pagestyle{fancy}
\pagenumbering{arabic}
\setcounter{page}{3}

% Rozdziały pracy
\newpage
% Tymczasowo wyłączone (wersja częściowa: literatura + projekt); pliki zawierają jeszcze tekst wzorcowy formatki.
% \section*{Streszczenie}
%\addcontentsline{toc}{chapter}{Streszczenie}
\label{chap:Streszczenie}
%\addcontentsline{toc}{chapter}{Streszczenie}
%\markboth{STRESZCZENIE}{STRESZCZENIE}

Lorem ipsum dolor sit amet, consectetur adipiscing elit. Sed viverra vitae justo non interdum. Nullam aliquam felis vel tellus efficitur, eget ultrices elit varius. Maecenas euismod ante sed tristique. Suspendisse vehicula efficitur ante, ac condimentum dui congue ut. Fusce in velit eu erat elementum dignissim eget non tortor. Proin sit amet est sit amet ex vehicula laoreet. Vivamus ut leo vel nisi volutpat facilisis. Praesent feugiat risus eu est gravida, id convallis nunc scelerisque. Nullam in lorem dui. Sed sit amet bibendum mauris, vel hendrerit est. Integer cursus odio eget metus sagittis, eu mattis ante vulputate. Sed hen-drerit, erat a aliquet bibendum, quam libero ultricies tortor, in sagittis ligula elit sed libero.
Pellentesque id hendrerit est. Sed gravida eros in justo accumsan, non sceler-isque dolor tincidunt. Vestibulum id nunc euismod, fermentum odio nec, tincid-unt arcu. Nunc eget leo nec eros euismod scelerisque. Nunc facilisis purus id neque laoreet interdum. Duis lacinia vestibulum mi, ut viverra est volutpat eget. Nullam aliquam nibh eu dui elementum scelerisque. Fusce efficitur libero a tel-lus convallis, in venenatis libero pellentesque. Suspendisse potenti. Fusce lacin-ia sapien ac odio tincidunt, nec feugiat tellus bibendum. Proin aliquam ex eget libero fermentum, in bibendum urna ultricies. Vivamus ut pellentesque libero. Aenean tincidunt dolor et quam hendrerit, vel ultricies ligula faucibus. Vivamus a elit eu dui bibendum fermentum id id neque. Nulla eget tortor euismod, tristique neque in, congue arcu.
Praesent tristique eget quam non dignissim. Integer placerat enim et erat con-sectetur, ut scelerisque nulla interdum. Pellentesque habitant morbi tristique senectus et netus et malesuada fames ac turpis egestas. Sed at lectus non orci dictum consectetur. Sed sagittis, eros at eleifend dignissim, metus arcu sollici-tudin libero, non egestas justo libero vel mi. Sed a urna arcu. Fusce sodales eros vel nulla fringilla posuere. 

\bigskip

\textbf{Słowa kluczowe}: \quad nie \quad więcej\quad niż \quad pięć

\textbf{Keywords}: \quad not \quad more \quad than \quad five


\newpage
% Tymczasowo wyłączone (wersja częściowa: literatura + projekt); pliki zawierają jeszcze tekst wzorcowy formatki.
% \section*{Wstęp}
\addcontentsline{toc}{section}{Wstęp}
\label{chap:Wstęp}
\markboth{WSTĘP}{WSTĘP}
Wstęp w pracy inżynierskiej wprowadza czytelnika w kontekst tematu pracy, zapoznając go z ogólnym tłem, które otacza badany problem lub projekt. We wstępie powinno się znaleźć krótkie przedstawienie obszaru, w którym została przeprowadzona praca, wskazanie aktualności problemu oraz podkreślenie jego znaczenia. Uświadom czytelnikowi, dlaczego omawiane zagadnienie jest ważne w kontekście inżynierii i technologii. Możesz podać krótką historię badanego zagadnienia, główne osiągnięcia w tym obszarze oraz wspomnieć o ewentualnych lukach w wiedzy lub problemach, które nadal pozostają nierozwiązane. Uzasadnij, dlaczego zdecydowałeś się na ten konkretny temat.

Jakie wydarzenia, trendy lub potrzeby rynku skłoniły Cię do podjęcia się tego zagadnienia. Przedstaw w skrócie używane metody badawcze lub podejście do analizy. Możesz wspomnieć, jakie narzędzia, technologie czy metody zostały użyte w pracy, bez wchodzenia w szczegóły (te zostaną przedstawione w osobnym rozdziale). Opisz, jak zorganizowana jest Twoja praca. Przedstaw krótko tematykę każdego rozdziału, aby czytelnik wiedział, czego się spodziewać. Pamiętaj, że wstęp ma na celu zainteresowanie czytelnika i przygotowanie go do głównego tematu pracy. Ma budować kontekst i stawiać Twoje badanie na tle innych prac w tej dziedzinie. Warto też dodać, że niektóre prace inżynierskie mogą mieć specyficzne wymagania dotyczące struktury wstępu, więc warto dostosować się do wytycznych wydziału czy promotora. Pamiętaj że powinieneś we wstępie dobrze nastroić recenzenta ;). Wstęp najczęściej piszemy już po skończeniu wszystkich rozdziałów pracy, zwykle ma około strony, chociaż w żadnej mierze nie oznacza to, że tak być musi. Bardzo rozsądnie jest poprosić o przeczytanie \textbf{całej} pracy kogoś, kto nie jest głęboko zaangażowany w jej realizację, być może nie jest nawet informatykiem (mama, tata, siostra, brat, ...). Ktoś taki skupi się na tekście i wskaże fragmenty niejasne. Jeżeli przyjmujecie taką radę, w przypadku wątpliwości, proszę nie tłumaczyć czytelnikowi dlaczego nie ma racji, ale poprawić tekst tak, by sam zrozumiał intencje autora.

\newpage
\section*{Cel i zakres pracy}
\addcontentsline{toc}{section}{Cel i zakres pracy}
\label{chap:cel}

\subsection*{Uzasadnienie podjęcia tematu}
Przeglądarki agentowe, czyli narzędzia wykorzystujące duże modele językowe do samodzielnego poruszania się po stronach internetowych i~wykonywania na nich operacji w~imieniu użytkownika, w~ciągu ostatnich dwóch lat przeszły z~etapu prototypów badawczych do etapu produktów udostępnianych szerokiemu gronu odbiorców. Ich naturalnym obszarem zastosowań jest automatyzacja powtarzalnych czynności biurowych, takich jak uzupełnianie formularzy, wyszukiwanie rekordów czy zmiana statusów w~panelach administracyjnych. Skuteczność tych narzędzi pozostaje jednak ograniczona, a~przyczyną niepowodzeń jest nie tylko niedoskonałość modeli, lecz również budowa samych aplikacji, projektowanych wyłącznie z~myślą o~obsłudze przez człowieka.

Istniejące zestawy testów porównawczych i~środowiska badawcze, omówione w~rozdziale~\ref{chap:pierwszy}, służą przede wszystkim do porównywania agentów i~modeli na niezmiennym zbiorze stron. Brakuje natomiast narzędzi pozwalających ocenić drugą stronę interakcji -- aplikację internetową -- oraz zmierzyć, w~jakim stopniu konkretne decyzje dotyczące budowy interfejsu wpływają na powodzenie zadań realizowanych przez agenta. Wypełnienie tej luki ma znaczenie praktyczne: twórcy aplikacji biurowych nie dysponują obecnie empirycznie zweryfikowanymi zaleceniami, które pozwalałyby przygotować interfejs do obsługi przez agentów, w~sposób analogiczny do wytycznych WCAG dotyczących dostępności dla osób z~niepełnosprawnościami.

\subsection*{Cel pracy}
Celem pracy jest \textbf{zaprojektowanie oraz implementacja środowiska badawczego i~zestawu testów porównawczych umożliwiających rzetelny pomiar skuteczności agenta przeglądarkowego Google Antigravity w~zadaniach realizowanych w~otwartoźródłowych aplikacjach biurowych}, z~obiektywną weryfikacją poprawności wykonania zadań na podstawie stanu danych po stronie serwera.

Opracowane środowisko posłuży następnie do zbadania, w~jakim stopniu modyfikacje struktury i~semantyki interfejsu aplikacji -- przy niezmienionym agencie, danych i~treści poleceń -- wpływają na poprawność realizacji zadań, oraz do sformułowania zaleceń dotyczących projektowania aplikacji internetowych przyjaznych agentom.

Cel główny jest niezależny od wielkości różnic, jakie zostaną zaobserwowane pomiędzy wariantami interfejsu: rzetelne środowisko pomiarowe stanowi wartość samą w~sobie, a~wynik eksperymentu -- zarówno potwierdzający, jak i~obalający postawione hipotezy -- jest wnioskiem z~badania, a~nie warunkiem jego powodzenia.

\subsection*{Cele szczegółowe}
Realizacja celu głównego wymaga osiągnięcia następujących celów szczegółowych:
\begin{enumerate}[label=C\arabic*.]
    \item Wybór, uruchomienie w~kontrolowanym środowisku lokalnym oraz przystosowanie co najmniej jednej otwartoźródłowej aplikacji biurowej typu panel administracyjny, z~zachowaniem jej oryginalnego interfejsu jako wariantu bazowego.
    \item Opracowanie co najmniej jednego wariantu interfejsu przyjaznego agentom, korzystającego z~tej samej logiki i~tych samych danych co wariant bazowy.
    \item Opracowanie zestawu scenariuszy testowych o~zróżnicowanym poziomie trudności, odzwierciedlających typowe operacje biurowe.
    \item Implementacja mechanizmu weryfikacji stanu danych, sprawdzającego zarówno wystąpienie oczekiwanych zmian, jak i~brak niezamierzonych skutków ubocznych.
    \item Opracowanie powtarzalnej procedury badawczej dla agenta Google Antigravity, obejmującej rejestrację konfiguracji agenta, przebiegu zadania i~jego wyniku.
    \item Przeprowadzenie eksperymentów porównujących skuteczność agenta w~poszczególnych wariantach interfejsu i~przy użyciu różnych modeli językowych.
    \item Analiza błędów agenta oraz sformułowanie empirycznie uzasadnionych zaleceń projektowych.
\end{enumerate}

\subsection*{Kryteria oceny}
Ocena realizacji celu opiera się na wielkościach mierzalnych, wyznaczanych automatycznie na podstawie stanu danych, a~nie na deklaracjach agenta ani na subiektywnej ocenie badacza. Podstawową miarą jest wskaźnik skuteczności realizacji zadań (TSR), czyli odsetek przebiegów zakończonych poprawną zmianą stanu danych bez skutków ubocznych, a~wpływ modyfikacji interfejsu wyrażany jest jako różnica tego wskaźnika pomiędzy wariantami. Miary uzupełniające obejmują czas wykonania zadania, liczbę kroków agenta, odsetek niezamierzonych skutków ubocznych, odsetek fałszywych deklaracji sukcesu, liczbę interwencji badacza oraz stabilność wyników przy wielokrotnym powtórzeniu tego samego scenariusza. Szczegółowe definicje miar przedstawiono w~podrozdziale~\ref{sec:procedura}.

\subsection*{Zakres pracy}
Zakres pracy obejmuje:
\begin{enumerate}
    \item \textbf{Przegląd literatury i~analizę istniejących rozwiązań} -- omówienie współczesnych agentów przeglądarkowych opartych na dużych modelach językowych i~modelach multimodalnych, sposobów reprezentacji strony internetowej wykorzystywanych przez agentów (drzewo DOM, drzewo dostępności, zrzuty ekranu), związku dostępności cyfrowej z~dostępnością dla agentów oraz krytyczną analizę istniejących zestawów testów porównawczych i~metod weryfikacji poprawności wykonania zadań (rozdział~\ref{chap:pierwszy}).
    \item \textbf{Projekt środowiska badawczego} -- sformułowanie wymagań, architektury, kryteriów wyboru aplikacji, zasad konstruowania wariantów interfejsu, formatu scenariuszy, mechanizmu weryfikacji stanu, procedury badawczej, miar i~hipotez (rozdział~\ref{chap:drugi}).
    \item \textbf{Wybór, adaptację i~modyfikację aplikacji biurowych} -- uruchomienie wybranych otwartoźródłowych paneli administracyjnych z~lokalną warstwą danych oraz implementację wariantów interfejsu różniących się m.in.~semantyką elementów HTML, nazewnictwem elementów interaktywnych, sposobem etykietowania pól formularzy, użyciem atrybutów ARIA oraz konstrukcją złożonych komponentów.
    \item \textbf{Przygotowanie scenariuszy testowych i~mechanizmu weryfikacji} -- opracowanie zadań na trzech poziomach trudności (operacje elementarne, złożone oraz wieloetapowe i~warunkowe) wraz z~warunkami poprawności i~listą dozwolonych zmian stanu.
    \item \textbf{Integrację z~agentem przeglądarkowym Google Antigravity} -- konfigurację agenta wywoływanego poleceniem \texttt{/browser}, przygotowanie procedury uruchamiania przebiegów oraz rejestracji ich wyników, z~możliwie wysokim stopniem automatyzacji.
    \item \textbf{Przeprowadzenie eksperymentów i~analizę wyników} -- wielokrotne wykonanie scenariuszy w~wariancie bazowym i~w~wariantach zmodyfikowanych, przy użyciu co najmniej dwóch modeli językowych o~różnych możliwościach, klasyfikację błędów agenta oraz ocenę istotności statystycznej zaobserwowanych różnic.
    \item \textbf{Opracowanie zaleceń projektowych i~wniosków} -- sformułowanie zestawu zaleceń dotyczących projektowania interfejsów przyjaznych agentom, odniesienie ich do kryteriów WCAG, wskazanie ograniczeń zastosowanego podejścia oraz kierunków dalszych badań.
\end{enumerate}

\subsection*{Ograniczenia zakresu}
W~celu zachowania spójności i~wykonalności pracy z~jej zakresu świadomie wyłączono następujące zagadnienia:
\begin{itemize}
    \item budowę własnego agenta przeglądarkowego, w~tym potoku przetwarzania złożonego z~narzędzia sterującego przeglądarką, reprezentacji strony i~modelu językowego -- badany jest wyłącznie gotowy agent dostępny w~środowisku Google Antigravity, traktowany jako czarna skrzynka;
    \item porównywanie różnych przeglądarek agentowych lub produktów komercyjnych -- jedyną zmienianą cechą po stronie agenta jest wykorzystywany model językowy;
    \item tworzenie aplikacji biurowych od podstaw -- wykorzystywane są istniejące projekty otwartoźródłowe, co zapewnia realizm interfejsu;
    \item trenowanie i~dostrajanie modeli językowych oraz budowę własnych systemów rozpoznawania obrazu;
    \item badanie stron dostępnych publicznie w~sieci Internet, których stanu nie można kontrolować;
    \item opracowanie formalnego standardu o~skali porównywalnej z~WCAG -- formułowane zalecenia mają charakter praktyczny i~ograniczają się do czynników zbadanych w~eksperymentach.
\end{itemize}


\newpage
\section{Przegląd literatury i istniejących rozwiązań}
\label{chap:pierwszy}
Rozdział stanowi techniczne wprowadzenie do zagadnienia skuteczności agentów przeglądarkowych. Omówiono w~nim podstawowe pojęcia, kierunki rozwoju agentów opartych na dużych modelach językowych, sposoby reprezentowania stron internetowych wykorzystywane przez takie systemy, a~także istniejące zestawy testów porównawczych i~metody weryfikacji poprawności wykonania zadań. Szczególną uwagę poświęcono związkom pomiędzy dostępnością cyfrową aplikacji a~skutecznością agentów, ponieważ to właśnie ten obszar wyznacza miejsce rozwiązania proponowanego w~niniejszej pracy. Rozdział kończy się zestawieniem porównawczym omówionych rozwiązań oraz wskazaniem luki, której wypełnienie stanowi przedmiot pracy.

% =====================================================================
\subsection{Podstawowe pojęcia}
\label{sec:pojecia}

\textbf{Agent} to w~klasycznym ujęciu sztucznej inteligencji system, który postrzega otoczenie za pomocą sensorów i~oddziałuje na nie za pomocą efektorów, dążąc do realizacji określonego celu \citep{russell2021aima}. W~kontekście niniejszej pracy otoczeniem agenta jest przeglądarka internetowa wraz z~wyświetlaną w~niej aplikacją, ,,sensorami'' są różne reprezentacje stanu strony (opisane w~podrozdziale~\ref{sec:reprezentacje}), a~,,efektorami'' -- akcje takie jak kliknięcie, wpisanie tekstu, przewinięcie strony czy przejście pod wskazany adres.

\textbf{Duży model językowy} (ang. \textit{large language model}, LLM) to model statystyczny oparty zazwyczaj na architekturze transformatora, wytrenowany na bardzo dużym zbiorze tekstów, zdolny do generowania tekstu oraz -- po odpowiednim dostrojeniu -- do wykonywania poleceń i~wywoływania narzędzi. \textbf{Model multimodalny} rozszerza tę zdolność o~przetwarzanie innych modalności, w~szczególności obrazów, co pozwala mu bezpośrednio interpretować zrzuty ekranu. Przegląd architektur agentów budowanych na bazie dużych modeli językowych przedstawiono w~pracy \cite{wang2024survey}.

\textbf{Agent przeglądarkowy} (nazywany również agentem internetowym, ang. \textit{web agent}) to system wykorzystujący model językowy lub multimodalny do samodzielnego wykonywania zadań w~przeglądarce internetowej na podstawie polecenia sformułowanego w~języku naturalnym. Działanie agenta przebiega w~pętli: obserwacja stanu strony, wnioskowanie o~kolejnym kroku, wykonanie akcji i~ponowna obserwacja. Schemat ten, łączący naprzemienne rozumowanie i~działanie, został spopularyzowany przez metodę ReAct \citep{yao2023react}. \textbf{Przeglądarką agentową} nazywa się natomiast gotowy produkt (aplikację lub usługę), w~którym agent przeglądarkowy jest dostępny dla użytkownika końcowego.

\textbf{Zadanie} rozumiane jest w~pracy jako polecenie w~języku naturalnym (np. ,,zmień status zamówienia o~numerze X na dostarczone''), któremu odpowiada jednoznacznie określony docelowy stan aplikacji. \textbf{Wskaźnik skuteczności realizacji zadań} (ang. \textit{task success rate}, TSR) to odsetek zadań zakończonych osiągnięciem stanu docelowego. Jest to podstawowa miara stosowana w~większości omawianych dalej zestawów testów porównawczych.

\textbf{Dostępność cyfrowa} to cecha aplikacji internetowej oznaczająca możliwość korzystania z~niej przez osoby o~różnych potrzebach, w~tym przez osoby z~niepełnosprawnościami posługujące się technologiami wspomagającymi, takimi jak czytniki ekranu. Podstawowym standardem w~tym zakresie są wytyczne WCAG \citep{w3c2023wcag22}. W~pracy przyjęto analogiczne pojęcie \textbf{dostępności dla agentów}, oznaczające cechy aplikacji ułatwiające jej poprawną obsługę przez agentów przeglądarkowych.

% =====================================================================
\subsection{Agenci przeglądarkowi oparci na dużych modelach językowych}
\label{sec:agenci}

\subsubsection{Rozwój podejść badawczych}
Pierwsze systematyczne prace nad agentami obsługującymi interfejsy internetowe dotyczyły uczenia ze wzmocnieniem w~uproszczonych, syntetycznych środowiskach. Platforma World of Bits wraz z~zestawem MiniWoB \citep{shi2017wob}, rozszerzonym następnie do MiniWoB++ \citep{liu2018miniwob}, zawierała kilkadziesiąt, a~następnie ponad sto niewielkich zadań, takich jak kliknięcie przycisku o~określonej etykiecie, wypełnienie prostego formularza czy wybór daty w~kalendarzu. Środowisko WebShop \citep{yao2022webshop} przeniosło ten problem na grunt bardziej realistyczny -- symulowanego sklepu internetowego zawierającego ponad milion produktów -- w~którym agent musiał wyszukać i~zakupić towar spełniający opis słowny.

Upowszechnienie dużych modeli językowych zmieniło sposób budowy agentów: zamiast trenować wyspecjalizowaną strategię dla każdego środowiska, wykorzystuje się model ogólnego przeznaczenia, któremu w~poleceniu przekazuje się opis zadania, tekstową reprezentację strony oraz listę dozwolonych akcji. Zbiór danych Mind2Web \citep{deng2023mind2web}, obejmujący ponad dwa tysiące zadań z~137 rzeczywistych witryn, pokazał, że już na etapie wyboru właściwego elementu strony z~bardzo długiego dokumentu HTML modele popełniają liczne błędy. Wraz z~pojawieniem się modeli multimodalnych powstały agenci operujący bezpośrednio na zrzutach ekranu, m.in. WebVoyager \citep{he2024webvoyager} oraz SeeAct \citep{zheng2024seeact}. Druga z~tych prac wykazała, że głównym ograniczeniem agentów wizualnych nie jest rozumienie zadania, lecz tzw. ugruntowanie (ang. \textit{grounding}), czyli powiązanie zamierzonej akcji z~konkretnym elementem interfejsu.

Istotnym wnioskiem płynącym z~nowszych badań jest to, że sposób przedstawienia strony agentowi ma wpływ na skuteczność porównywalny z~wyborem samego modelu. Autorzy rozwiązania AgentOccam \citep{yang2025agentoccam} uzyskali znaczącą poprawę wyników w~środowisku WebArena wyłącznie poprzez uproszczenie przestrzeni obserwacji i~akcji, bez zmiany modelu językowego. Spostrzeżenie to ma bezpośrednie znaczenie dla niniejszej pracy: jeżeli uproszczenie reprezentacji strony po stronie agenta poprawia jego działanie, to uzasadnione jest pytanie, czy podobny efekt można uzyskać, zmieniając samą aplikację.

\subsubsection{Komercyjne przeglądarki agentowe}
Od 2024 roku agenci przeglądarkowi przestali być wyłącznie przedmiotem badań i~zaczęli pojawiać się w~produktach komercyjnych. Firma Anthropic udostępniła mechanizm sterowania komputerem (ang. \textit{computer use}), w~którym model na podstawie zrzutów ekranu generuje ruchy kursora i~naciśnięcia klawiszy \citep{anthropic2024computeruse}. OpenAI zaprezentowało model Computer-Using Agent, stanowiący podstawę usługi Operator \citep{openai2025cua}. Google DeepMind rozwija projekt Mariner \citep{google2024mariner}, a~w~październiku 2025 roku udostępniło wyspecjalizowany model Gemini 2.5 Computer Use, przeznaczony do sterowania przeglądarką na podstawie zrzutów ekranu \citep{google2025computeruse}.

Wspólną cechą tych rozwiązań jest to, że użytkownik nie ma wpływu na wewnętrzną budowę agenta -- otrzymuje gotowy system, którego sposób postrzegania strony, dobór modelu i~strategia działania pozostają w~dużej mierze nieujawnione. Z~perspektywy twórcy aplikacji internetowej jest to sytuacja analogiczna do tej, w~jakiej znajduje się on wobec przeglądarek internetowych czy czytników ekranu: nie kontroluje programu klienckiego, a~jedynie strukturę dostarczanego dokumentu.

\subsubsection{Agent przeglądarkowy w~środowisku Google Antigravity}
\label{sec:antigravity}
Google Antigravity to środowisko programistyczne nowej generacji, zaprezentowane przez Google w~listopadzie 2025 roku wraz z~rodziną modeli Gemini~3 \citep{google2025antigravityblog}. W~odróżnieniu od klasycznych asystentów programisty, środowisko to opiera się na agentach zdolnych do samodzielnego planowania i~wykonywania wieloetapowych zadań w~trzech obszarach: edytorze kodu, terminalu oraz przeglądarce internetowej.

Obsługę przeglądarki realizuje wyspecjalizowany podagent przeglądarkowy, wywoływany m.in. poleceniem \texttt{/browser} \citep{google2025antigravitydocs}. Zgodnie z~dokumentacją producenta podagent ten steruje lokalną instancją przeglądarki Google Chrome uruchamianą z~odrębnym profilem użytkownika, działa we własnym kontekście rozmowy, odseparowanym od głównego agenta, a~przebieg jego pracy jest rejestrowany w~postaci zrzutów ekranu i~nagrań wideo. Środowisko umożliwia również ograniczenie zbioru adresów, które agent może odwiedzać. Użytkownik może wybrać model językowy, z~którego korzysta agent, co pozwala porównywać skuteczność tego samego narzędzia przy różnych modelach.
% TODO (autor): zweryfikować powyższy opis z aktualną wersją dokumentacji Antigravity
% (https://antigravity.google/docs/ide/browser) i uzupełnić o numer wersji używanej w badaniach.

Szczegóły wewnętrznego działania agenta -- w~szczególności to, w~jakim stopniu wykorzystuje on zrzuty ekranu, a~w~jakim strukturę dokumentu lub drzewo dostępności -- nie są w~pełni udokumentowane publicznie. Z~tego powodu w~pracy agent Antigravity traktowany jest jako system typu ,,czarna skrzynka'': przedmiotem pomiaru jest wyłącznie jego zachowanie i~skuteczność w~zależności od konstrukcji obsługiwanej aplikacji. Wybór tego narzędzia uzasadniają trzy przesłanki: jest to aktualny, rozwijany produkt dużego dostawcy, udostępnia możliwość zmiany modelu bez zmiany pozostałych elementów systemu, a~sposób jego użycia odpowiada rzeczywistemu scenariuszowi, w~którym użytkownik końcowy zleca agentowi pracę w~aplikacji biurowej.

% =====================================================================
\subsection{Reprezentacje strony internetowej wykorzystywane przez agentów}
\label{sec:reprezentacje}
Agent przeglądarkowy nie ,,widzi'' strony tak jak człowiek; otrzymuje jedną lub kilka jej reprezentacji. Od wyboru reprezentacji zależy zarówno to, jakie informacje są dla agenta dostępne, jak i~to, które właściwości aplikacji mogą wpływać na jego skuteczność.

\subsubsection{Obiektowy model dokumentu}
Obiektowy model dokumentu (ang. \textit{Document Object Model}, DOM) to drzewiasta reprezentacja dokumentu HTML udostępniana przez przeglądarkę skryptom i~narzędziom zewnętrznym \citep{whatwg2025dom}. Przekazanie agentowi kodu HTML lub serializacji drzewa DOM zapewnia pełny dostęp do struktury strony, w~tym do atrybutów elementów, nazw klas i~identyfikatorów. Zasadniczym ograniczeniem tego podejścia jest rozmiar danych: według raportu WebAIM Million z~2025 roku przeciętna strona główna popularnej witryny zawiera ponad 1250 elementów \citep{webaim2025million}, a~w~rozbudowanych aplikacjach jednostronicowych liczba ta jest wielokrotnie większa. Drzewo DOM zawiera ponadto wiele elementów pozbawionych znaczenia dla użytkownika (kontenery układu, elementy dekoracyjne), a~nazwy klas generowane przez biblioteki stylów nie niosą informacji semantycznej. Dlatego agenci rzadko korzystają z~pełnego drzewa DOM i~zazwyczaj stosują jego filtrowanie lub uproszczenie \citep{deng2023mind2web}.

\subsubsection{Drzewo dostępności}
Drzewo dostępności (ang. \textit{accessibility tree}) to wyprowadzana przez przeglądarkę z~drzewa DOM reprezentacja przeznaczona pierwotnie dla technologii wspomagających. Każdy węzeł tego drzewa opisany jest przez \textbf{rolę} (np. przycisk, pole tekstowe, łącze, wiersz tabeli), \textbf{dostępną nazwę} (ang. \textit{accessible name}), \textbf{stan} (np. zaznaczony, rozwinięty, nieaktywny) oraz \textbf{wartość}. Role i~stany wynikają z~semantyki elementów HTML \citep{whatwg2025html} lub z~atrybutów WAI-ARIA \citep{w3c2023aria12}, a~sposób wyznaczania dostępnej nazwy -- na podstawie treści elementu, powiązanej etykiety, atrybutów \texttt{aria-label} i~\texttt{aria-labelledby} czy atrybutu \texttt{title} -- określa odrębna specyfikacja \citep{w3c2024accname}.

Drzewo dostępności jest znacznie zwięźlejsze niż DOM i~zawiera dokładnie te informacje, które są potrzebne do obsługi interfejsu: jakie elementy są interaktywne, do czego służą i~w~jakim są stanie. Z~tego powodu stanowi ono podstawową reprezentację obserwacji w~środowisku WebArena \citep{zhou2024webarena} i~w~wielu innych rozwiązaniach. Jakość tej reprezentacji jest jednak w~całości zależna od twórcy aplikacji. Element \texttt{<div>} z~dołączoną obsługą zdarzenia kliknięcia nie posiada w~drzewie dostępności roli przycisku, a~przycisk zawierający wyłącznie ikonę graficzną nie posiada nazwy. Dla agenta korzystającego z~drzewa dostępności takie elementy są zatem niewidoczne lub nieodróżnialne od siebie, mimo że dla człowieka mogą być w~pełni czytelne.

\subsubsection{Zrzuty ekranu i~podejście wizualne}
Agenci wizualni otrzymują obraz wyrenderowanej strony, co pozwala im interpretować układ, ikony i~elementy graficzne w~sposób zbliżony do człowieka. Rozwiązanie to jest niezależne od jakości kodu HTML, lecz wprowadza inne trudności: konieczność precyzyjnego wskazania współrzędnych elementu, ograniczoną rozdzielczość, problem elementów znajdujących się poza widocznym obszarem oraz trudność w~odczytaniu stanów słabo wyróżnionych wizualnie. W~celu ograniczenia problemu ugruntowania stosuje się technikę oznaczania zbioru elementów (ang. \textit{Set-of-Mark}), polegającą na nałożeniu na zrzut ekranu numerowanych ramek wokół elementów interaktywnych \citep{yang2023som}. Ramki te wyznaczane są jednak zazwyczaj na podstawie drzewa DOM lub drzewa dostępności, co oznacza, że również agenci określani jako wizualni pośrednio korzystają ze struktury dokumentu \citep{he2024webvoyager, koh2024visualwebarena}.

\subsubsection{Reprezentacje hybrydowe}
Współczesne rozwiązania coraz częściej łączą powyższe reprezentacje: zrzut ekranu służy do zrozumienia kontekstu wizualnego, a~drzewo dostępności lub uproszczony DOM -- do jednoznacznej identyfikacji elementów i~wykonywania akcji. Taki model obserwacji oferuje m.in. platforma BrowserGym \citep{chezelles2025browsergym}. Zestawienie cech poszczególnych reprezentacji przedstawiono w~tabeli~\ref{tab:reprezentacje}.

\begin{table}[H]
\centering
\caption{Porównanie reprezentacji strony internetowej wykorzystywanych przez agentów (opracowanie własne na podstawie \citep{zhou2024webarena, zheng2024seeact, chezelles2025browsergym})}
\label{tab:reprezentacje}
\footnotesize
\begin{tabularx}{\textwidth}{>{\raggedright\arraybackslash}p{2.4cm}>{\raggedright\arraybackslash}X>{\raggedright\arraybackslash}X>{\raggedright\arraybackslash}X}
\toprule
\textbf{Reprezentacja} & \textbf{Zalety} & \textbf{Ograniczenia} & \textbf{Właściwości aplikacji wpływające na agenta} \\
\midrule
Drzewo DOM / kod HTML & pełna informacja o~strukturze i~atrybutach & bardzo duży rozmiar, wiele elementów bez znaczenia, brak informacji o~wyglądzie & struktura dokumentu, atrybuty, liczba zagnieżdżeń \\
\addlinespace
Drzewo dostępności & zwięzłość, jawne role, nazwy i~stany elementów & całkowita zależność od poprawności semantyki i~atrybutów ARIA & semantyka HTML, etykiety, atrybuty ARIA, komunikaty o~stanie \\
\addlinespace
Zrzut ekranu & niezależność od kodu, dostęp do informacji wizualnej & problem ugruntowania, ograniczony obszar widoczny, niejednoznaczne ikony & czytelność wizualna, rozmieszczenie elementów, widoczność stanów \\
\addlinespace
Hybrydowa & łączy kontekst wizualny z~jednoznaczną identyfikacją elementów & większy koszt obliczeniowy, możliwa niespójność reprezentacji & wszystkie powyższe \\
\bottomrule
\end{tabularx}
\end{table}

Z~punktu widzenia niniejszej pracy najistotniejsze jest spostrzeżenie zawarte w~ostatniej kolumnie tabeli: niezależnie od przyjętej reprezentacji istnieją właściwości aplikacji, które twórca może kontrolować i~które potencjalnie wpływają na skuteczność agenta. Ponieważ sposób postrzegania strony przez agenta Antigravity nie jest w~pełni znany, badanie wpływu tych właściwości wymaga pomiaru empirycznego.

% =====================================================================
\subsection{Dostępność cyfrowa a~dostępność dla agentów}
\label{sec:dostepnosc}

\subsubsection{Standardy dostępności}
Wytyczne WCAG 2.2 \citep{w3c2023wcag22} formułują wymagania dostępności w~oparciu o~cztery zasady: postrzegalność, funkcjonalność, zrozumiałość i~kompatybilność. Kilka kryteriów sukcesu tych wytycznych ma bezpośrednie odniesienie do problemów obserwowanych u~agentów przeglądarkowych:
\begin{itemize}
    \item kryterium 1.3.1 ,,Informacje i~relacje'' -- struktura i~powiązania między elementami (np. pole formularza i~jego etykieta, nagłówek tabeli i~komórka) muszą być określone programowo, a~nie wyłącznie wizualnie;
    \item kryterium 2.4.6 ,,Nagłówki i~etykiety'' oraz 3.3.2 ,,Etykiety lub instrukcje'' -- elementy powinny posiadać etykiety opisujące ich przeznaczenie;
    \item kryterium 4.1.2 ,,Nazwa, rola, wartość'' -- dla każdego komponentu interfejsu nazwa i~rola muszą być możliwe do ustalenia programowo, a~stany i~wartości -- do odczytania i~zmiany;
    \item kryterium 4.1.3 ,,Komunikaty o~stanie'' -- komunikaty o~powodzeniu operacji, błędach czy postępie muszą być udostępniane technologiom wspomagającym bez konieczności przenoszenia na nie fokusu.
\end{itemize}

Uzupełnieniem WCAG jest specyfikacja WAI-ARIA \citep{w3c2023aria12}, definiująca role, stany i~właściwości umożliwiające opisanie niestandardowych komponentów (np. list rozwijanych z~wyszukiwaniem, okien modalnych, zakładek). Specyfikacja ta oraz towarzyszące jej zalecenia podkreślają jednak, że w~pierwszej kolejności należy korzystać z~natywnych elementów HTML posiadających wbudowaną semantykę, a~atrybuty ARIA stosować dopiero wtedy, gdy odpowiedni element natywny nie istnieje. Niepoprawne użycie ARIA może bowiem pogorszyć dostępność w~porównaniu z~jej brakiem.

Mimo istnienia standardów poziom dostępności stron internetowych pozostaje niski. Według raportu WebAIM Million z~2025 roku automatycznie wykrywalne naruszenia WCAG występowały na 94,8\% spośród miliona najpopularniejszych stron głównych, a~do najczęstszych należały m.in. brak tekstu alternatywnego dla obrazów oraz brak etykiet pól formularzy \citep{webaim2025million}. Znaczenie dostępności rośnie również z~przyczyn prawnych: dyrektywa Parlamentu Europejskiego i~Rady (UE) 2019/882, stosowana od czerwca 2025 roku, rozszerzyła wymagania dostępności na wiele produktów i~usług sektora prywatnego \citep{ue2019eaa}.

\subsubsection{Dostępność jako czynnik skuteczności narzędzi automatycznych}
Związek między dostępnością a~możliwością automatycznej obsługi interfejsu jest od dawna wykorzystywany w~testowaniu oprogramowania. Biblioteka Testing Library zaleca, aby testy wyszukiwały elementy przede wszystkim po roli i~dostępnej nazwie, a~dopiero w~dalszej kolejności po etykiecie, tekście zastępczym czy identyfikatorze testowym, uzasadniając to tym, że testy powinny jak najwierniej odzwierciedlać sposób, w~jaki aplikacja jest używana \citep{testinglibrary2025priority}. Podobne zalecenie zawiera dokumentacja narzędzia Playwright, które jako preferowany sposób lokalizowania elementów wskazuje metody oparte na roli i~dostępnej nazwie \citep{playwright2025locators}. Aplikacja poprawnie zbudowana pod względem dostępności jest zatem łatwiejsza do obsługi przez programy automatyczne, a~agent przeglądarkowy -- w~odróżnieniu od skryptu testowego -- musi samodzielnie, bez wcześniejszej znajomości aplikacji, odnaleźć odpowiednie elementy.

\subsubsection{Koncepcje interfejsów przeznaczonych dla agentów}
W~2025 roku pojawiły się prace postulujące odwrócenie dotychczasowej perspektywy badań. \cite{lu2025awi} w~pracy o~znamiennym tytule ,,Build the web for agents, not agents for the web'' argumentują, że zamiast budować coraz bardziej złożonych agentów dostosowanych do interfejsów projektowanych dla ludzi, należy projektować interfejsy dostosowane do agentów. Autorzy proponują koncepcję agentowego interfejsu internetowego (ang. \textit{Agentic Web Interface}, AWI) i~formułują zasady jego projektowania, takie jak bezpieczeństwo, efektywność i~standaryzacja. Praca ma jednak charakter stanowiska (ang. \textit{position paper}) i~nie zawiera badań empirycznych wpływu konkretnych decyzji projektowych na skuteczność istniejących agentów.

Równolegle w~ramach społeczności W3C rozwijana jest propozycja WebMCP \citep{w3c2025webmcp}, umożliwiająca stronie internetowej jawne udostępnienie agentowi zbioru narzędzi (funkcji opisanych w~języku naturalnym wraz ze schematem parametrów), dzięki którym agent może wykonywać operacje bez symulowania kliknięć. Propozycja ta nawiązuje do protokołu Model Context Protocol \citep{anthropic2024mcp}. Rozwiązania tego typu wymagają jednak wsparcia zarówno po stronie aplikacji, jak i~agenta, a~ich upowszechnienie jest kwestią przyszłości. W~okresie przejściowym agenci będą nadal obsługiwać aplikacje poprzez ich interfejs graficzny, dlatego zasadne jest badanie, w~jaki sposób konstrukcja tego interfejsu wpływa na skuteczność agentów.

% =====================================================================
\subsection{Istniejące zestawy testów porównawczych i~środowiska badawcze}
\label{sec:benchmarki}
W~niniejszym podrozdziale omówiono najważniejsze zestawy testów porównawczych służące do oceny agentów przeglądarkowych. Dla każdego z~nich wskazano zalety oraz ograniczenia z~perspektywy celu niniejszej pracy.

\subsubsection{Środowiska syntetyczne: MiniWoB++ i~WebShop}
MiniWoB++ \citep{liu2018miniwob} zawiera ponad sto niewielkich zadań realizowanych na sztucznie przygotowanych stronach, z~automatycznie wyznaczaną nagrodą. Zaletą zestawu jest pełna kontrola nad środowiskiem i~możliwość wyizolowania pojedynczych typów interakcji (np. obsługi kalendarza czy listy rozwijanej). Strony te są jednak bardzo uproszczone i~nie odzwierciedlają złożoności rzeczywistych aplikacji. WebShop \citep{yao2022webshop} oferuje bardziej realistyczny scenariusz sklepu internetowego, lecz ogranicza się do jednego typu zadania (wyszukanie i~zakup produktu).

\subsubsection{Mind2Web}
Mind2Web \citep{deng2023mind2web} to zbiór zadań zebranych na rzeczywistych witrynach wraz z~zapisanymi sekwencjami akcji wykonanych przez ludzi. Ocena odbywa się na podstawie wcześniej zapisanych migawek stron, bez interakcji z~działającą aplikacją -- poprzez porównanie akcji przewidzianych przez model z~zapisem wzorcowym na poziomie pojedynczych kroków. Zaletą jest duża różnorodność witryn; ograniczeniem -- brak obserwacji rzeczywistych skutków akcji agenta oraz to, że ocena kroków nie uwzględnia alternatywnych, równie poprawnych ścieżek wykonania zadania. Kontynuacja tej pracy, Online-Mind2Web \citep{xue2025illusion}, wykazała, że w~ocenie na rzeczywistych witrynach skuteczność wielu agentów jest wyraźnie niższa niż wynikałoby to z~wcześniej publikowanych rezultatów, co autorzy określili mianem ,,iluzji postępu''.

\subsubsection{WebArena i~VisualWebArena}
WebArena \citep{zhou2024webarena} jest pod względem ogólnej idei rozwiązaniem najbliższym niniejszej pracy. Środowisko składa się z~kilku samodzielnie uruchamianych, w~pełni funkcjonalnych aplikacji internetowych (sklepu internetowego, panelu administracyjnego systemu sklepu Magento, forum dyskusyjnego, systemu GitLab oraz serwisu mapowego), uzupełnionych o~narzędzia pomocnicze. Zestaw obejmuje 812 zadań, a~ich poprawność weryfikowana jest funkcjonalnie: poprzez porównanie odpowiedzi agenta z~odpowiedzią wzorcową lub przez programowe sprawdzenie stanu aplikacji (np. zapytanie do bazy danych lub odczyt zawartości strony) po zakończeniu zadania. W~pierwotnej publikacji najlepszy agent oparty na modelu GPT-4 osiągnął skuteczność 14,41\%, podczas gdy ludzie -- 78,24\%. VisualWebArena \citep{koh2024visualwebarena} rozszerza to podejście o~zadania wymagające rozumienia treści wizualnych.

Z~perspektywy niniejszej pracy WebArena dostarcza dwóch kluczowych wzorców: wykorzystania realistycznych aplikacji uruchamianych w~kontrolowanym środowisku oraz weryfikacji opartej na stanie systemu. Jednocześnie różni się od niej w~sposób zasadniczy. W~WebArena agent jest budowany przez badacza według schematu: narzędzie automatyzacji przeglądarki $\rightarrow$ reprezentacja strony $\rightarrow$ model językowy $\rightarrow$ akcja $\rightarrow$ narzędzie automatyzacji przeglądarki. Przedmiotem oceny jest model lub architektura agenta, natomiast aplikacje pozostają niezmienne. Środowisko nie przewiduje tworzenia wariantów interfejsu ani nie formułuje zaleceń dla twórców aplikacji.

\subsubsection{WorkArena i~platforma BrowserGym}
WorkArena \citep{drouin2024workarena} to zestaw zadań realizowanych na platformie ServiceNow, czyli w~systemie klasy przedsiębiorstwa wykorzystywanym do obsługi procesów biurowych (zgłoszeń, katalogów usług, list i~formularzy). Jest to więc rozwiązanie najbliższe niniejszej pracy pod względem dziedziny zadań. Autorzy udostępnili również platformę BrowserGym \citep{chezelles2025browsergym}, ujednolicającą sposób definiowania środowisk, obserwacji i~akcji dla agentów przeglądarkowych. Podobnie jak WebArena, rozwiązanie to służy do porównywania agentów i~modeli, a~badana aplikacja (komercyjna platforma ServiceNow) nie może być dowolnie modyfikowana przez badacza.

\subsubsection{WebVoyager}
WebVoyager \citep{he2024webvoyager} to agent multimodalny wraz z~zestawem 643 zadań realizowanych na 15 rzeczywistych, publicznie dostępnych witrynach. Zaletą jest realistyczny charakter interakcji i~pełne, wieloetapowe zadania. Ograniczeniem -- brak kontroli nad testowanymi witrynami: ich treść i~wygląd zmieniają się w~czasie, co obniża powtarzalność wyników. Ponieważ stanu zewnętrznych serwisów nie można odczytać, poprawność zadań oceniana jest przez ludzi lub automatycznie przez model multimodalny analizujący zrzuty ekranu, co wprowadza element subiektywności.

\subsubsection{WebSuite}
WebSuite \citep{li2024websuite} jest zestawem diagnostycznym, którego celem nie jest wyznaczenie łącznej skuteczności agenta, lecz ustalenie przyczyn jego niepowodzeń. Autorzy zaproponowali taksonomię akcji wykonywanych w~interfejsach internetowych i~przygotowali aplikację, w~której każde zadanie testuje pojedynczy typ akcji lub komponentu. Pozwala to wskazać konkretne elementy interfejsu, z~którymi agenci mają trudności. Praca ta stanowi ważną inspirację dla części eksperymentalnej niniejszej pracy, ponieważ łączy niepowodzenia agentów z~konstrukcją interfejsu. Jej ograniczeniem jest skupienie na izolowanych komponentach -- zadania nie odzwierciedlają kompletnych, wieloetapowych scenariuszy w~realistycznej aplikacji biznesowej.

\subsubsection{BrowseComp i~BrowseComp-Plus}
BrowseComp \citep{wei2025browsecomp} obejmuje 1266 trudnych pytań, na które odpowiedź wymaga wytrwałego przeszukiwania sieci, przy czym poprawność odpowiedzi jest łatwa do sprawdzenia. Zestaw ten koncentruje się na wyszukiwaniu informacji, a~nie na wykonywaniu operacji w~aplikacjach. BrowseComp-Plus \citep{chen2025browsecompplus} zastępuje otwartą sieć stałym korpusem około 100 tysięcy dokumentów, co zwiększa powtarzalność i~pozwala oddzielić wpływ systemu wyszukiwania od wpływu modelu. Rozwiązanie to jest wartościowe metodologicznie -- pokazuje znaczenie kontrolowanego środowiska dla rzetelności pomiaru -- lecz nie dotyczy interakcji z~interfejsem użytkownika.

\subsubsection{Inne środowiska}
Warto wspomnieć również o~środowisku OSWorld \citep{xie2024osworld}, w~którym agenci wykonują zadania w~pełnym systemie operacyjnym, a~poprawność weryfikowana jest przez dedykowane skrypty sprawdzające stan aplikacji i~plików. Zakres tego środowiska wykracza poza przeglądarkę internetową, jednak przyjęta w~nim metoda weryfikacji jest zbieżna z~metodą proponowaną w~niniejszej pracy.

% =====================================================================
\subsection{Metody weryfikacji poprawności wykonania zadań}
\label{sec:weryfikacja}
Na podstawie przeglądu literatury można wyróżnić cztery podstawowe metody oceny, czy agent poprawnie wykonał zadanie:
\begin{enumerate}
    \item \textbf{Porównanie odpowiedzi} -- stosowane w~zadaniach informacyjnych; odpowiedź agenta porównywana jest z~wzorcową w~sposób dokładny lub przybliżony (np. BrowseComp, część zadań WebArena). Metoda nie nadaje się do zadań polegających na zmianie stanu aplikacji.
    \item \textbf{Porównanie ze ścieżką wzorcową} -- akcje agenta porównywane są z~zapisem akcji wykonanych przez człowieka (Mind2Web). Metoda nie uwzględnia alternatywnych poprawnych ścieżek i~nie potwierdza, że zamierzony efekt rzeczywiście nastąpił.
    \item \textbf{Ocena przez człowieka lub przez model} -- człowiek lub model multimodalny ocenia przebieg zadania na podstawie zrzutów ekranu i~deklaracji agenta (WebVoyager, Online-Mind2Web). Metoda jest uniwersalna, lecz kosztowna lub podatna na błędy oceniającego modelu.
    \item \textbf{Programowa weryfikacja stanu} -- po zakończeniu zadania skrypt sprawdza rzeczywisty stan aplikacji, np. zawartość bazy danych (WebArena, WorkArena, OSWorld). Metoda jest obiektywna i~powtarzalna, lecz wymaga pełnej kontroli nad środowiskiem oraz przygotowania procedury sprawdzającej dla każdego zadania.
\end{enumerate}

Szczególnie istotny jest problem rozbieżności między deklaracją agenta a~rzeczywistym stanem systemu. Agent może zakomunikować wykonanie zadania, mimo że zmiana nie została zapisana (np. formularz nie został zatwierdzony lub operacja została odrzucona z~powodu niepoprawnych danych), albo może wykonać zadanie poprawnie, lecz dodatkowo zmodyfikować inne rekordy. Wiarygodną ocenę zapewnia zatem jedynie programowa weryfikacja stanu, uzupełniona o~sprawdzenie braku niezamierzonych zmian. To podejście przyjęto w~niniejszej pracy.

% =====================================================================
\subsection{Zestawienie istniejących rozwiązań i~luka badawcza}
\label{sec:luka}
Zestawienie omówionych rozwiązań pod kątem kryteriów istotnych dla celu pracy przedstawiono w~tabeli~\ref{tab:porownanie}.

\begin{table}[H]
\centering
\caption{Porównanie istniejących zestawów testów porównawczych z~rozwiązaniem proponowanym w~pracy (opracowanie własne)}
\label{tab:porownanie}
\scriptsize
\begin{tabularx}{\textwidth}{>{\raggedright\arraybackslash}p{2.0cm}>{\raggedright\arraybackslash}X>{\raggedright\arraybackslash}p{1.6cm}>{\raggedright\arraybackslash}X>{\raggedright\arraybackslash}p{1.9cm}>{\raggedright\arraybackslash}p{1.5cm}}
\toprule
\textbf{Rozwiązanie} & \textbf{Badane aplikacje} & \textbf{Kontrolowane środowisko} & \textbf{Metoda weryfikacji} & \textbf{Przedmiot oceny} & \textbf{Warianty interfejsu} \\
\midrule
MiniWoB++ & syntetyczne strony z~pojedynczymi zadaniami & tak & programowa (nagroda) & agent & nie \\
WebShop & symulowany sklep internetowy & tak & programowa (zgodność atrybutów) & agent & nie \\
Mind2Web & 137 rzeczywistych witryn (zapisane migawki) & częściowo & zgodność ze ścieżką wzorcową & model & nie \\
WebArena, VisualWebArena & samodzielnie uruchamiane aplikacje (sklep, panel administracyjny, forum, GitLab) & tak & odpowiedź i~programowa weryfikacja stanu & agent, model & nie \\
WorkArena & platforma ServiceNow (system biurowy) & tak & programowa weryfikacja stanu & agent, model & nie \\
WebVoyager & 15 rzeczywistych witryn & nie & ocena przez człowieka lub model & agent & nie \\
WebSuite & aplikacja diagnostyczna z~izolowanymi komponentami & tak & programowa dla pojedynczych akcji & przyczyny błędów agenta & nie \\
BrowseComp & otwarta sieć & nie & porównanie odpowiedzi & agent (wyszukiwanie) & nie \\
BrowseComp-Plus & stały korpus dokumentów & tak & porównanie odpowiedzi & agent i~system wyszukiwania & nie \\
\midrule
\textbf{Rozwiązanie proponowane} & \textbf{otwartoźródłowe panele administracyjne} & \textbf{tak} & \textbf{programowa weryfikacja stanu wraz z~kontrolą skutków ubocznych} & \textbf{aplikacja (interfejs)} & \textbf{tak} \\
\bottomrule
\end{tabularx}
\end{table}

Analiza literatury prowadzi do następujących wniosków:
\begin{enumerate}
    \item Istniejące zestawy testów porównawczych odpowiadają przede wszystkim na pytania o~to, który model lub agent jest skuteczniejszy oraz jak ulepszyć agenta, aby lepiej radził sobie z~istniejącymi stronami. Badane aplikacje traktowane są jako niezmienne tło eksperymentu.
    \item Rozwiązania o~największej rzetelności pomiaru (WebArena, WorkArena, OSWorld) wymagają budowy własnego agenta lub integracji z~określonym interfejsem programistycznym i~nie odzwierciedlają sytuacji użytkownika gotowej przeglądarki agentowej.
    \item Jedynym spośród omówionych rozwiązań łączącym niepowodzenia agentów z~konstrukcją interfejsu jest WebSuite, lecz dotyczy ono izolowanych komponentów, a~nie kompletnych scenariuszy w~realistycznych aplikacjach.
    \item Postulat projektowania interfejsów pod kątem agentów \citep{lu2025awi} nie został dotychczas poparty badaniami empirycznymi wskazującymi, które konkretne decyzje projektowe wpływają na skuteczność istniejących agentów i~w~jakim stopniu.
    \item Zbieżność kryteriów WCAG (w~szczególności 1.3.1, 4.1.2 i~4.1.3) z~typowymi przyczynami błędów agentów sugeruje, że praktyki dostępności cyfrowej mogą jednocześnie zwiększać dostępność aplikacji dla agentów. Hipoteza ta wymaga jednak weryfikacji, szczególnie w~odniesieniu do agentów multimodalnych, które teoretycznie mogłyby radzić sobie z~interfejsem wyłącznie na podstawie obrazu.
\end{enumerate}

Luką badawczą, której wypełnienie stanowi przedmiot niniejszej pracy, jest zatem brak środowiska i~metody pozwalających w~kontrolowany sposób mierzyć wpływ konstrukcji interfejsu realistycznej aplikacji biurowej na skuteczność gotowego, komercyjnego agenta przeglądarkowego. Rozwiązanie wypełniające tę lukę powinno łączyć: realistyczne aplikacje biurowe, gotowego agenta traktowanego jako czarna skrzynka, powtarzalne scenariusze testowe, programową weryfikację stanu, kontrolowane warianty tego samego interfejsu oraz możliwość sformułowania na podstawie wyników zaleceń dla twórców aplikacji. Projekt takiego rozwiązania przedstawiono w~rozdziale~\ref{chap:drugi}.


\newpage
\section{Projekt środowiska badawczego}
\label{chap:drugi}
Rozdział przedstawia projekt środowiska badawczego służącego do pomiaru skuteczności agenta przeglądarkowego Google Antigravity w~zadaniach realizowanych w~aplikacjach biurowych. Omówiono w~nim założenia projektowe wynikające z~analizy literatury, wymagania, architekturę systemu, kryteria wyboru badanych aplikacji, sposób konstruowania wariantów interfejsu, format scenariuszy testowych, mechanizm weryfikacji stanu oraz procedurę badawczą wraz z~metrykami i~hipotezami. Opis implementacji poszczególnych modułów zostanie przedstawiony w~dalszej części pracy.

% =====================================================================
\subsection{Założenia projektowe na tle istniejących rozwiązań}
\label{sec:zalozenia}
Podstawowym założeniem projektu jest zmiana perspektywy w~stosunku do rozwiązań omówionych w~rozdziale~\ref{chap:pierwszy}. Typowe podejście polega na ulepszaniu agenta tak, aby poradził sobie z~istniejącą stroną; w~niniejszej pracy agent pozostaje niezmienny, a~zmienną eksperymentu stanowi aplikacja. Pytanie badawcze brzmi zatem: \emph{w~jaki sposób zmienić aplikację, aby istniejący agent wykonywał w~niej zadania skuteczniej i~bardziej niezawodnie?}

Z~tego założenia wynikają trzy decyzje o~charakterze ograniczającym. Po pierwsze, w~pracy nie jest budowany własny agent przeglądarkowy -- wykorzystywany jest gotowy agent dostępny w~środowisku Google Antigravity, traktowany jako czarna skrzynka. Po drugie, nie są tworzone od podstaw własne aplikacje testowe -- wykorzystywane są istniejące, otwartoźródłowe panele administracyjne, co zapewnia realizm interfejsu i~ogranicza ryzyko nieświadomego dopasowania aplikacji do agenta. Po trzecie, praca nie ma na celu porównywania różnych przeglądarek agentowych ani trenowania modeli językowych.

Projekt czerpie jednocześnie z~rozwiązań istniejących, adaptując ich sprawdzone elementy do nowego celu. Powiązanie inspiracji z~literatury z~decyzjami projektowymi przedstawiono w~tabeli~\ref{tab:inspiracje}.

\begin{table}[H]
\centering
\caption{Elementy istniejących rozwiązań wykorzystane w~projekcie (opracowanie własne)}
\label{tab:inspiracje}
\footnotesize
\begin{tabularx}{\textwidth}{>{\raggedright\arraybackslash}p{2.6cm}>{\raggedright\arraybackslash}X>{\raggedright\arraybackslash}X}
\toprule
\textbf{Źródło} & \textbf{Przejęty element} & \textbf{Adaptacja w~projekcie} \\
\midrule
WebArena \citep{zhou2024webarena} & realistyczne aplikacje uruchamiane lokalnie; weryfikacja stanu po zakończeniu zadania & aplikacje biurowe typu panel administracyjny z~lokalną warstwą danych; weryfikacja rozszerzona o~kontrolę skutków ubocznych \\
\addlinespace
WorkArena \citep{drouin2024workarena} & zadania z~dziedziny pracy biurowej (listy, formularze, zmiany statusów) & scenariusze obejmujące typowe operacje biurowe w~systemach otwartoźródłowych, które można modyfikować \\
\addlinespace
WebSuite \citep{li2024websuite} & powiązanie niepowodzeń agenta z~konkretnymi typami komponentów & klasyfikacja błędów agenta i~przypisanie ich do elementów interfejsu; warianty komponentów jako zmienna eksperymentu \\
\addlinespace
BrowseComp-Plus \citep{chen2025browsecompplus} & kontrola środowiska w~celu oddzielenia wpływu poszczególnych czynników & identyczne dane początkowe i~polecenia dla wszystkich wariantów; jedyną zmienną jest interfejs \\
\addlinespace
WCAG, WAI-ARIA \citep{w3c2023wcag22, w3c2023aria12} & kryteria nazwy, roli, wartości, relacji i~komunikatów o~stanie & źródło hipotez dotyczących modyfikacji interfejsu i~punkt odniesienia dla formułowanych zaleceń \\
\addlinespace
AWI \citep{lu2025awi} & postulat dostosowania interfejsu do agentów & empiryczna weryfikacja postulatu w~ramach istniejących technologii (HTML, ARIA), bez zmiany agenta \\
\bottomrule
\end{tabularx}
\end{table}

% =====================================================================
\subsection{Wymagania}
\label{sec:wymagania}
Na podstawie celu pracy oraz wniosków z~analizy literatury sformułowano wymagania funkcjonalne (WF) i~niefunkcjonalne (WN) stawiane środowisku badawczemu.

\textbf{Wymagania funkcjonalne:}
\begin{itemize}
    \item[WF-01] Środowisko umożliwia uruchomienie co najmniej jednej otwartoźródłowej aplikacji biurowej w~co najmniej dwóch wariantach interfejsu, korzystających z~tej samej logiki i~tych samych danych.
    \item[WF-02] Środowisko umożliwia przywrócenie stanu danych aplikacji do jednoznacznie określonego stanu początkowego przed każdym przebiegiem zadania.
    \item[WF-03] Scenariusze testowe są definiowane w~jednolitym, czytelnym formacie zawierającym polecenie dla agenta, poziom trudności, stan początkowy oraz warunki poprawności.
    \item[WF-04] Po zakończeniu przebiegu środowisko automatycznie sprawdza stan danych i~określa, czy zadanie zostało wykonane poprawnie oraz czy nie wystąpiły niezamierzone zmiany.
    \item[WF-05] Wynik każdego przebiegu jest zapisywany wraz z~informacją o~scenariuszu, wariancie interfejsu, konfiguracji agenta (w~tym użytym modelu), czasie trwania oraz szczegółowym wynikiem weryfikacji.
    \item[WF-06] Środowisko umożliwia zestawienie wyników w~przekroju: konfiguracja agenta $\times$ scenariusz $\times$ wariant interfejsu.
\end{itemize}

\textbf{Wymagania niefunkcjonalne:}
\begin{itemize}
    \item[WN-01] Powtarzalność -- uruchomienie środowiska na innym komputerze prowadzi do identycznego stanu początkowego aplikacji i~danych.
    \item[WN-02] Izolacja zmiennej eksperymentu -- warianty interfejsu różnią się wyłącznie badanymi właściwościami; dane, logika i~polecenia pozostają identyczne.
    \item[WN-03] Nieujawnianie wariantu -- agent nie może na podstawie adresu, tytułu strony ani treści polecenia ustalić, który wariant interfejsu obsługuje.
    \item[WN-04] Obiektywność -- ocena poprawności nie zależy od deklaracji agenta ani od subiektywnej oceny badacza.
    \item[WN-05] Rozszerzalność -- dodanie nowego scenariusza, wariantu lub aplikacji nie wymaga zmian w~pozostałych modułach.
    \item[WN-06] Wierność wariantu bazowego -- wariant bazowy odpowiada oryginalnej aplikacji i~nie jest celowo pogarszany.
\end{itemize}

% =====================================================================
\subsection{Architektura rozwiązania}
\label{sec:architektura}
Architekturę środowiska badawczego przedstawiono na rysunku~\ref{fig:architektura}. System ma budowę modułową; jego elementy można podzielić na dwie grupy: elementy, z~którymi agent wchodzi w~interakcję (prawa część rysunku), oraz elementy wspierające prowadzenie badania, niewidoczne dla agenta (lewa część rysunku).

\begin{figure}[H]
\centering
\begin{tikzpicture}[
    node distance=7mm and 24mm,
    box/.style={draw, rounded corners=2pt, align=center, text width=46mm, minimum height=11mm, font=\footnotesize},
    agent/.style={box, fill=gray!15},
    arr/.style={-{Stealth[length=2.2mm]}, thick},
    lbl/.style={font=\scriptsize, align=center, fill=white, inner sep=1pt}
]
% prawa kolumna - elementy widoczne dla agenta
\node[agent] (ag) {Agent przeglądarkowy\\Google Antigravity\\(\texttt{/browser})\\+ wybrany model językowy};
\node[box, below=of ag] (chr) {Przeglądarka Chrome\\(odrębny profil)};
\node[box, below=of chr] (app) {Aplikacja biurowa\\wariant bazowy /\\wariant przyjazny agentom};
\node[box, below=of app] (api) {Lokalny serwer danych (API)};
\node[box, below=of api] (db) {Baza danych\\(stan aplikacji)};
% lewa kolumna - elementy badawcze
\node[box, left=of ag] (sc) {Moduł scenariuszy\\(definicje zadań)};
\node[box, left=of chr] (st) {Moduł sterujący\\przebiegiem badania};
\node[box, left=of api] (wer) {Moduł weryfikacji stanu};
\node[box, left=of db] (wyn) {Repozytorium wyników\\i~moduł analizy};

\draw[arr] (ag) -- node[lbl, right]{akcje} (chr);
\draw[arr] (chr) -- (app);
\draw[arr] (app) -- node[lbl, right]{żądania HTTP} (api);
\draw[arr] (api) -- (db);
\draw[arr] (sc) -- (st);
% polecenie i przywrocenie stanu prowadzone pionowo w odstepie miedzy kolumnami
\coordinate (gap) at ($(st.east)!0.5!(chr.west)$);
\draw[thick] (st.east) -- (gap);
\draw[arr] (gap) |- node[lbl, pos=0.25]{polecenie} (ag.west);
\draw[arr] (gap) |- node[lbl, pos=0.25]{przywrócenie\\stanu} ([yshift=2mm]api.west);
\draw[arr] ([yshift=-2mm]api.west) -- node[lbl, below]{odczyt stanu} ([yshift=-2mm]wer.east);
\draw[arr] (st) -- node[lbl, left]{wyzwolenie} (wer);
\draw[arr] (wer) -- (wyn);

\begin{scope}[on background layer]
\node[draw, dashed, inner sep=3mm, fit=(ag)(db), label={[font=\scriptsize]above:elementy widoczne dla agenta}] {};
\node[draw, dashed, inner sep=3mm, fit=(sc)(wyn), label={[font=\scriptsize]above:elementy wspierające badanie}] {};
\end{scope}
\end{tikzpicture}
\caption{Architektura środowiska badawczego (opracowanie własne)}
\label{fig:architektura}
\end{figure}

Poszczególne moduły pełnią następujące funkcje:
\begin{itemize}
    \item \textbf{Agent przeglądarkowy Google Antigravity} -- badany system; otrzymuje polecenie w~języku naturalnym i~wykonuje je, sterując przeglądarką. Jego konfiguracja (wersja środowiska, model językowy) jest rejestrowana dla każdego przebiegu.
    \item \textbf{Aplikacja biurowa} -- otwartoźródłowy panel administracyjny uruchamiany lokalnie, w~każdym z~badanych wariantów interfejsu (podrozdziały~\ref{sec:aplikacje} i~\ref{sec:warianty}).
    \item \textbf{Lokalny serwer danych i~baza danych} -- warstwa danych współdzielona przez wszystkie warianty interfejsu; stanowi źródło prawdy o~stanie aplikacji. Udostępnia dodatkowo niewidoczny dla agenta punkt końcowy służący do przywracania stanu początkowego i~odczytu pełnego stanu danych.
    \item \textbf{Moduł scenariuszy} -- zbiór definicji zadań testowych (podrozdział~\ref{sec:scenariusze}).
    \item \textbf{Moduł sterujący} -- realizuje procedurę pojedynczego przebiegu: przywraca stan początkowy, przekazuje polecenie agentowi, rejestruje czas, a~po zakończeniu zadania uruchamia weryfikację.
    \item \textbf{Moduł weryfikacji stanu} -- porównuje stan danych po zakończeniu przebiegu ze stanem oczekiwanym (podrozdział~\ref{sec:weryfikacja-projekt}).
    \item \textbf{Repozytorium wyników i~moduł analizy} -- przechowuje wyniki przebiegów i~umożliwia ich zestawianie oraz porównywanie wariantów (podrozdział~\ref{sec:procedura}).
\end{itemize}

Środowisko uruchomieniowe aplikacji i~warstwy danych zostanie zamknięte w~kontenerach (np. z~wykorzystaniem narzędzia Docker), co zapewni spełnienie wymagania powtarzalności WN-01. Agent Antigravity działa natomiast na komputerze badacza i~łączy się z~aplikacją poprzez adres lokalny.

% =====================================================================
\subsection{Wybór aplikacji biurowych}
\label{sec:aplikacje}
\subsubsection{Kryteria wyboru}
Badane aplikacje powinny spełniać następujące kryteria:
\begin{enumerate}
    \item \textbf{Realizm} -- aplikacja odzwierciedla rzeczywisty system biurowy i~zawiera typowe elementy: nawigację, listy z~wyszukiwaniem, filtrowaniem i~stronicowaniem, formularze tworzenia i~edycji rekordów, zmiany statusów oraz okna modalne.
    \item \textbf{Otwarta licencja} -- kod źródłowy jest dostępny na licencji pozwalającej na modyfikację na potrzeby badań.
    \item \textbf{Możliwość lokalnego uruchomienia} -- aplikacja działa bez zależności od zewnętrznych usług, których stanu nie można kontrolować.
    \item \textbf{Dostęp do stanu danych} -- możliwe jest programowe odczytanie i~przywrócenie stanu danych.
    \item \textbf{Modyfikowalność interfejsu} -- komponenty interfejsu mogą być zastępowane bez ingerencji w~logikę aplikacji.
    \item \textbf{Zróżnicowanie} -- wybrane aplikacje korzystają z~różnych bibliotek komponentów, co pozwala ocenić, czy obserwowane zależności mają charakter ogólny.
\end{enumerate}

\subsubsection{Analizowane aplikacje}
Wstępnie przeanalizowano sześć otwartoźródłowych aplikacji typu panel administracyjny. Lista ta ma charakter otwarty i~może zostać rozszerzona. Zestawienie ich cech przedstawiono w~tabeli~\ref{tab:aplikacje}.

\begin{table}[H]
\centering
\caption{Porównanie analizowanych otwartoźródłowych aplikacji biurowych (opracowanie własne na podstawie dokumentacji projektów)}
\label{tab:aplikacje}
\scriptsize
\begin{tabularx}{\textwidth}{>{\raggedright\arraybackslash}p{2.3cm}>{\raggedright\arraybackslash}p{2.6cm}>{\raggedright\arraybackslash}p{2.0cm}>{\raggedright\arraybackslash}X>{\raggedright\arraybackslash}p{1.3cm}}
\toprule
\textbf{Aplikacja} & \textbf{Dziedzina} & \textbf{Biblioteka komponentów} & \textbf{Warstwa danych i~uwagi} & \textbf{Licencja} \\
\midrule
react-admin, przykład \textit{Posters Galore} \citep{reactadmin} & obsługa sklepu: zamówienia, klienci, produkty, recenzje & Material UI & dane generowane w~przeglądarce; wymienny moduł dostępu do danych umożliwia podłączenie lokalnego serwera & MIT \\
\addlinespace
refine, przykład \textit{finefoods-antd} \citep{refine} & obsługa dostaw żywności: zamówienia, produkty, kurierzy, sklepy & Ant Design & domyślnie zdalny serwer demonstracyjny; wymienny moduł dostępu do danych & MIT \\
\addlinespace
Twenty \citep{twenty} & system zarządzania relacjami z~klientami (CRM) & własna biblioteka komponentów & pełna część serwerowa z~bazą PostgreSQL; wysoki realizm, duża złożoność uruchomienia i~modyfikacji & AGPL-3.0 \\
\addlinespace
Saleor Dashboard \citep{saleor} & panel administracyjny sklepu internetowego & Macaw UI & wymaga uruchomienia systemu Saleor (interfejs GraphQL, PostgreSQL); duża złożoność & BSD-3-Clause \\
\addlinespace
shadcn-admin \citep{shadcnadmin} & szablon panelu: zadania, użytkownicy, aplikacje & shadcn/ui, Radix UI & dane statyczne; wymaga dodania warstwy danych & MIT \\
\addlinespace
TailAdmin \citep{tailadmin} & szablon panelu z~wykresami i~tabelami & Tailwind CSS & głównie widoki statyczne, ograniczona logika operacji na danych & MIT \\
\bottomrule
\end{tabularx}
\end{table}
% TODO (autor): potwierdzić licencje i stan projektów w dniu rozpoczęcia implementacji.

\subsubsection{Wybór aplikacji do badań}
Na podstawie przyjętych kryteriów w~pierwszej kolejności do badań przewidziano aplikacje \textbf{react-admin (Posters Galore)} oraz \textbf{refine (finefoods-antd)}. Obie zawierają pełny zestaw operacji tworzenia, odczytu, edycji i~usuwania danych na kilku powiązanych ze sobą zasobach, a~ich architektura przewiduje wymienny moduł dostępu do danych (ang. \textit{data provider}). Pozwala to zastąpić dane generowane w~przeglądarce lub zdalny serwer demonstracyjny lokalnym serwerem z~bazą danych, bez zmian w~warstwie prezentacji. Aplikacje te korzystają z~dwóch różnych, popularnych bibliotek komponentów (Material UI i~Ant Design), co umożliwi ocenę, czy obserwowane zależności nie są specyficzne dla jednej biblioteki. System Twenty, jako aplikacja o~najwyższym stopniu realizmu, rozważany jest jako rozszerzenie badań w~dalszym etapie. Szablony shadcn-admin i~TailAdmin, ze względu na brak rzeczywistej warstwy danych, mogą posłużyć co najwyżej jako źródło przykładów komponentów.
% TODO (autor): ostateczny wybór aplikacji do potwierdzenia z promotorem.

% =====================================================================
\subsection{Warianty interfejsu}
\label{sec:warianty}
\subsubsection{Zasada konstruowania wariantów}
Kluczowym elementem projektu jest możliwość uruchomienia tej samej aplikacji w~kilku wariantach interfejsu przy zachowaniu identycznej logiki i~danych:
\begin{itemize}
    \item \textbf{wariant bazowy} -- oryginalny interfejs aplikacji otwartoźródłowej, bez celowych modyfikacji (wymaganie WN-06); jedyną zmianą jest podłączenie lokalnej warstwy danych;
    \item \textbf{wariant przyjazny agentom} -- ta sama funkcjonalność, w~której wybrane komponenty zmodyfikowano zgodnie z~hipotezami badawczymi;
    \item opcjonalnie \textbf{warianty pośrednie} -- zawierające pojedynczą modyfikację, co pozwala określić wpływ każdego czynnika z~osobna.
\end{itemize}

Warianty implementowane są w~postaci alternatywnych wersji komponentów wybieranych na etapie budowania aplikacji. Każdy wariant udostępniany jest pod odrębnym portem lokalnym o~neutralnym oznaczeniu, a~tytuł strony i~pozostałe widoczne elementy nie zawierają informacji o~wariancie. Rozwiązanie to spełnia wymaganie WN-03: gdyby wariant był wskazywany np. parametrem w~adresie strony, agent mógłby go odczytać, co stanowiłoby potencjalne źródło zakłócenia eksperymentu.

Lista modyfikacji wprowadzanych w~wariancie przyjaznym agentom nie jest ustalana w~całości z~góry. Zgodnie z~procedurą opisaną w~podrozdziale~\ref{sec:procedura} zostanie ona uzupełniona na podstawie analizy błędów popełnianych przez agenta w~wariancie bazowym. Punktem wyjścia są czynniki wskazane w~literaturze oraz kryteria WCAG, zestawione w~tabeli~\ref{tab:czynniki}.

\begin{table}[H]
\centering
\caption{Badane czynniki i~przykładowe modyfikacje interfejsu (opracowanie własne na podstawie \citep{w3c2023wcag22, w3c2023aria12, li2024websuite})}
\label{tab:czynniki}
\scriptsize
\begin{tabularx}{\textwidth}{>{\raggedright\arraybackslash}p{2.2cm}>{\raggedright\arraybackslash}X>{\raggedright\arraybackslash}X>{\raggedright\arraybackslash}p{1.3cm}>{\raggedright\arraybackslash}p{1.0cm}}
\toprule
\textbf{Czynnik} & \textbf{Przykład w~wariancie bazowym} & \textbf{Modyfikacja w~wariancie przyjaznym agentom} & \textbf{Kryterium WCAG} & \textbf{Hipoteza} \\
\midrule
Semantyka elementów akcji & element \texttt{div} z~obsługą kliknięcia & natywny element \texttt{button} lub \texttt{a} & 4.1.2 & H1 \\
\addlinespace
Nazwy elementów akcji & przycisk z~samą ikoną; ogólne etykiety ,,OK'', ,,Zapisz'' & dostępna nazwa opisująca operację, np. ,,Zapisz zamówienie'', ,,Usuń klienta Jan Kowalski'' & 2.4.6, 4.1.2 & H2 \\
\addlinespace
Etykiety pól formularza & tekst zastępczy zamiast etykiety; etykieta niepowiązana z~polem & element \texttt{label} powiązany z~polem; grupowanie pól (\texttt{fieldset}) & 1.3.1, 3.3.2 & H2, H3 \\
\addlinespace
Złożone komponenty & niestandardowa lista rozwijana z~wyszukiwaniem, kalendarz, przełącznik & element natywny (\texttt{select}, \texttt{input type="date"}) lub komponent zgodny z~wzorcami WAI-ARIA & 4.1.2 & H4 \\
\addlinespace
Komunikaty o~stanie & powiadomienie znikające po kilku sekundach; zmiana sygnalizowana wyłącznie kolorem & trwały komunikat z~rolą \texttt{status} lub \texttt{alert}; jawny tekst stanu & 4.1.3, 1.4.1 & H4 \\
\addlinespace
Nawigacja & zwinięte menu z~samymi ikonami; zagnieżdżone menu rozwijane & widoczne menu tekstowe w~elemencie \texttt{nav}, oznaczenie bieżącej strony & 2.4.x & H2 \\
\addlinespace
Tabele danych & siatka zbudowana z~elementów \texttt{div} & element \texttt{table} z~nagłówkami \texttt{th} lub siatka zgodna z~rolą \texttt{grid} & 1.3.1 & H1 \\
\bottomrule
\end{tabularx}
\end{table}

Przykład pojedynczej modyfikacji przedstawiono na listingu~\ref{lst:wariant}. Oba fragmenty mogą zostać wyświetlone w~sposób wizualnie identyczny, jednak w~drzewie dostępności pierwszy z~nich nie jest widoczny jako element interaktywny i~nie posiada nazwy, natomiast drugi występuje jako przycisk o~nazwie ,,Usuń zamówienie nr 1024''.

\begin{lstlisting}[language=HTML, caption={Przycisk usuwania w~wariancie bazowym (u~góry) i~w~wariancie przyjaznym agentom (u~dołu)}, label=lst:wariant, basicstyle=\ttfamily\footnotesize]
<!-- wariant bazowy -->
<div class="row-action" onclick="deleteOrder(1024)">
  <svg class="icon-trash">...</svg>
</div>

<!-- wariant przyjazny agentom -->
<button type="button" onclick="deleteOrder(1024)"
        aria-label="Usuń zamówienie nr 1024">
  <svg class="icon-trash" aria-hidden="true">...</svg>
</button>
\end{lstlisting}

% =====================================================================
\subsection{Scenariusze testowe}
\label{sec:scenariusze}
Scenariusz testowy jest jawną i~kompletną definicją zadania. Każdy scenariusz zawiera: identyfikator, nazwę, aplikację, poziom trudności, identyfikator stanu początkowego, polecenie dla agenta w~języku naturalnym oraz listę warunków poprawności. Polecenie jest identyczne dla wszystkich wariantów interfejsu i~nie zawiera wskazówek dotyczących sposobu wykonania zadania (np. nazw przycisków czy kolejności kroków), aby nie faworyzować żadnego z~wariantów. Przykładową definicję scenariusza przedstawiono na listingu~\ref{lst:scenariusz}.

\begin{lstlisting}[language={}, caption={Przykładowa definicja scenariusza testowego}, label=lst:scenariusz, basicstyle=\ttfamily\footnotesize]
id: RA-07
nazwa: Zmiana statusu zamówienia
aplikacja: react-admin-posters
poziom: 2
stan_poczatkowy: dane-01
polecenie: >
  W panelu administracyjnym dostępnym pod adresem
  http://localhost:8101 znajdź zamówienie o numerze
  referencyjnym ABC123 i zmień jego status na
  "dostarczone".
warunki:
  - zasob: commands
    filtr: { reference: "ABC123" }
    oczekiwane: { status: "delivered" }
dozwolone_zmiany:
  - zasob: commands
    filtr: { reference: "ABC123" }
    pola: [status]
\end{lstlisting}

Scenariusze podzielono na trzy poziomy trudności:
\begin{enumerate}
    \item \textbf{Poziom 1 -- operacje elementarne}: pojedyncza operacja wymagająca niewielkiej liczby kroków, np. odnalezienie rekordu i~odczyt wartości, zmiana jednego pola.
    \item \textbf{Poziom 2 -- operacje złożone}: operacje wymagające obsługi złożonych komponentów lub formularzy z~wieloma polami, np. utworzenie nowego klienta, zmiana statusu z~użyciem listy rozwijanej, ustawienie kilku filtrów.
    \item \textbf{Poziom 3 -- operacje wieloetapowe i~warunkowe}: sekwencje kilku operacji na różnych zasobach lub zadania wymagające odczytania informacji i~podjęcia na jej podstawie decyzji, np. ,,znajdź klienta z~największą liczbą zamówień w~bieżącym miesiącu i~nadaj mu status VIP''.
\end{enumerate}
Zakłada się przygotowanie około 15--20 scenariuszy dla każdej badanej aplikacji, w~proporcjach zapewniających reprezentację wszystkich poziomów trudności oraz wszystkich czynników z~tabeli~\ref{tab:czynniki}.

% =====================================================================
\subsection{Mechanizm weryfikacji stanu}
\label{sec:weryfikacja-projekt}
Zgodnie z~wnioskami z~podrozdziału~\ref{sec:weryfikacja} ocena poprawności opiera się wyłącznie na programowej weryfikacji stanu danych, a~nie na deklaracjach agenta. Pojedyncza weryfikacja przebiega następująco:
\begin{enumerate}
    \item Przed rozpoczęciem przebiegu warstwa danych jest przywracana do stanu początkowego wskazanego w~scenariuszu, a~następnie zapisywana jest migawka pełnego stanu $S_0$.
    \item Po zakończeniu przebiegu odczytywany jest stan końcowy $S_1$.
    \item Sprawdzane są \textbf{warunki poprawności}: dla każdego warunku wyszukiwane są rekordy spełniające filtr i~porównywane wartości wskazanych pól z~wartościami oczekiwanymi. Obsługiwane są również warunki dotyczące utworzenia lub usunięcia rekordu oraz liczby rekordów.
    \item Wyznaczany jest \textbf{zbiór różnic} pomiędzy $S_0$ i~$S_1$; każda różnica nieobjęta listą dozwolonych zmian traktowana jest jako niezamierzony skutek uboczny.
    \item Przebieg uznaje się za zakończony sukcesem wtedy i~tylko wtedy, gdy spełnione są wszystkie warunki poprawności i~nie wystąpiły niezamierzone skutki uboczne.
\end{enumerate}

Kontrola skutków ubocznych stanowi rozszerzenie podejścia stosowanego w~środowisku WebArena, w~którym sprawdza się zasadniczo jedynie wystąpienie oczekiwanej zmiany. W~zastosowaniach biurowych przypadkowa modyfikacja innego rekordu (np. zmiana statusu niewłaściwego zamówienia obok właściwego) jest błędem równie poważnym jak niewykonanie zadania, dlatego nie może być pominięta w~ocenie.

Niezależnie od wyniku weryfikacji rejestrowana jest deklaracja agenta dotycząca wykonania zadania. Zestawienie deklaracji z~wynikiem weryfikacji pozwala wyznaczyć odsetek przypadków, w~których agent błędnie uznał zadanie za wykonane -- zjawisko szczególnie istotne z~punktu widzenia zaufania do agentów w~zastosowaniach biurowych.

% =====================================================================
\subsection{Procedura badawcza i~metryki}
\label{sec:procedura}
\subsubsection{Przebieg pojedynczego zadania}
Pojedynczy przebieg zadania obejmuje następujące kroki:
\begin{enumerate}
    \item przywrócenie stanu początkowego danych i~zapis migawki $S_0$,
    \item uruchomienie nowej, niezależnej sesji agenta Antigravity z~wybranym modelem, aby wyeliminować wpływ wcześniejszych rozmów,
    \item przekazanie agentowi polecenia ze scenariusza za pomocą polecenia \texttt{/browser},
    \item oczekiwanie na zakończenie pracy agenta lub przekroczenie ustalonego limitu czasu,
    \item zapis deklaracji agenta, czasu trwania oraz dostępnych zapisów przebiegu (zrzutów ekranu, nagrań, zapisu rozmowy),
    \item uruchomienie weryfikacji stanu i~zapis wyniku.
\end{enumerate}
W~trakcie przebiegu badacz nie ingeruje w~pracę agenta. Jeżeli agent zwróci się z~prośbą o~doprecyzowanie lub potwierdzenie, odpowiedź udzielana jest według ustalonego z~góry schematu, a~fakt ten rejestrowany jest jako interwencja. Stopień automatyzacji kroków 2--5 zależy od możliwości udostępnianych przez środowisko Antigravity (np. interfejsu wiersza poleceń lub zestawu narzędzi programistycznych dla języka Python); w~przypadku ich niewystarczalności kroki te wykonywane będą przez badacza według ustalonej instrukcji, a~pozostałe kroki pozostaną zautomatyzowane.
% TODO (autor): ustalić, czy uruchamianie agenta da się w pełni zautomatyzować (agy CLI / Antigravity SDK).

\subsubsection{Cykl eksperymentalny}
Badanie ma charakter iteracyjny, co przedstawiono na rysunku~\ref{fig:cykl}. Po wykonaniu wszystkich scenariuszy w~wariancie bazowym przeprowadzana jest analiza błędów agenta, na podstawie której formułowane są modyfikacje interfejsu. Następnie te same scenariusze wykonywane są w~wariancie zmodyfikowanym, a~wyniki porównywane są z~wynikami wariantu bazowego.

\begin{figure}[H]
\centering
\begin{tikzpicture}[
    node distance=14mm and 9mm,
    etap/.style={draw, rounded corners=2pt, align=center, text width=36mm, minimum height=13mm, font=\footnotesize},
    arr/.style={-{Stealth[length=2.2mm]}, thick}
]
\node[etap] (s1) {Wariant bazowy};
\node[etap, right=of s1] (s2) {Seria przebiegów\\wszystkich scenariuszy};
\node[etap, right=of s2] (s3) {Analiza i~klasyfikacja\\błędów agenta};
\node[etap, below=of s3] (s4) {Hipotezy\\i~modyfikacje interfejsu};
\node[etap, left=of s4] (s5) {Wariant zmodyfikowany\\i~ponowna seria\\przebiegów};
\node[etap, left=of s5] (s6) {Porównanie wyników\\i~zalecenia projektowe};
\draw[arr] (s1) -- (s2);
\draw[arr] (s2) -- (s3);
\draw[arr] (s3) -- (s4);
\draw[arr] (s4) -- (s5);
\draw[arr] (s5) -- (s6);
\draw[arr, dashed] (s5.north) -- node[font=\scriptsize, right]{kolejna iteracja} (s2.south);
\end{tikzpicture}
\caption{Iteracyjny cykl eksperymentalny (opracowanie własne)}
\label{fig:cykl}
\end{figure}

\subsubsection{Metryki}
Podstawową metryką jest wskaźnik skuteczności realizacji zadań, wyznaczany dla danej konfiguracji agenta $a$ i~wariantu interfejsu $v$ jako:
\begin{equation}
    \mathrm{TSR}(a, v) = \frac{1}{N \cdot K} \sum_{i=1}^{N} \sum_{k=1}^{K} s_{i,k}(a, v),
    \label{eq:tsr}
\end{equation}
\myequations{Wskaźnik skuteczności realizacji zadań}
\noindent
gdzie $N$ oznacza liczbę scenariuszy, $K$ -- liczbę powtórzeń każdego scenariusza, a~$s_{i,k}(a, v) \in \{0, 1\}$ -- wynik weryfikacji $k$-tego przebiegu $i$-tego scenariusza (1 -- sukces, 0 -- niepowodzenie).

Wpływ modyfikacji interfejsu wyrażany jest jako różnica wskaźników skuteczności, podawana w~punktach procentowych:
\begin{equation}
    \Delta(a) = \mathrm{TSR}(a, v_{\mathrm{przyjazny}}) - \mathrm{TSR}(a, v_{\mathrm{bazowy}}).
    \label{eq:delta}
\end{equation}
\myequations{Różnica skuteczności pomiędzy wariantami interfejsu}

Ponieważ działanie agenta opartego na modelu językowym jest niedeterministyczne, każdy scenariusz wykonywany jest wielokrotnie (zakłada się $K \geq 3$), a~oprócz wartości wskaźnika podawany jest jego przedział ufności. Ze względu na to, że te same scenariusze wykonywane są w~obu wariantach, do oceny istotności statystycznej różnic zostanie zastosowany test dla prób zależnych (np. test McNemara).

Metryki uzupełniające obejmują:
\begin{itemize}
    \item czas wykonania zadania,
    \item liczbę kroków (akcji) agenta -- o~ile jest możliwa do ustalenia na podstawie zapisu przebiegu,
    \item odsetek przebiegów z~niezamierzonymi skutkami ubocznymi,
    \item odsetek fałszywych deklaracji sukcesu (agent zadeklarował wykonanie zadania, a~weryfikacja zakończyła się niepowodzeniem),
    \item liczbę interwencji badacza,
    \item stabilność wyników, rozumianą jako odsetek scenariuszy, w~których wszystkie powtórzenia zakończyły się tym samym wynikiem.
\end{itemize}

\subsubsection{Klasyfikacja błędów}
Każdy nieudany przebieg jest analizowany na podstawie zapisu przebiegu i~przypisywany do jednej z~kategorii błędów: (1)~nieodnalezienie elementu interfejsu, (2)~wybór niewłaściwego elementu, (3)~nieprawidłowa obsługa komponentu (np. listy rozwijanej, kalendarza, okna modalnego), (4)~błąd nawigacji, (5)~błędna interpretacja stanu aplikacji (np. nieodczytanie komunikatu o~błędzie), (6)~błąd zrozumienia lub planowania zadania, (7)~przekroczenie limitu czasu lub zapętlenie. Kategorie 1--5 są potencjalnie zależne od konstrukcji interfejsu i~stanowią podstawę formułowania modyfikacji w~kolejnej iteracji; kategoria 6 dotyczy przede wszystkim możliwości modelu.

\subsubsection{Model danych wyników}
Wyniki przechowywane są zgodnie z~modelem danych przedstawionym w~tabeli~\ref{tab:model}, umożliwiającym zestawianie wyników w~przekroju wymaganym przez WF-06.

\begin{table}[H]
\centering
\caption{Model danych wyników badania (opracowanie własne)}
\label{tab:model}
\footnotesize
\begin{tabularx}{\textwidth}{>{\raggedright\arraybackslash}p{3.4cm}>{\raggedright\arraybackslash}X}
\toprule
\textbf{Encja} & \textbf{Najważniejsze atrybuty} \\
\midrule
KonfiguracjaAgenta & identyfikator, wersja środowiska Antigravity, model językowy, ustawienia \\
Scenariusz & identyfikator, aplikacja, nazwa, poziom trudności, polecenie, stan początkowy \\
WariantInterfejsu & identyfikator, aplikacja, opis modyfikacji, badane czynniki \\
Przebieg & identyfikator, konfiguracja, scenariusz, wariant, numer powtórzenia, czas rozpoczęcia i~zakończenia, deklaracja agenta, liczba interwencji, wynik, kategoria błędu \\
WynikWeryfikacji & identyfikator przebiegu, warunek, wartość oczekiwana, wartość rzeczywista, wynik; wykryte skutki uboczne \\
\bottomrule
\end{tabularx}
\end{table}

% =====================================================================
\subsection{Hipotezy badawcze}
\label{sec:hipotezy}
Na podstawie analizy literatury sformułowano następujące hipotezy, które zostaną zweryfikowane z~wykorzystaniem zaprojektowanego środowiska:
\begin{itemize}
    \item[\textbf{H1}] Zastąpienie niestandardowych elementów interaktywnych semantycznymi elementami HTML o~identycznym wyglądzie zwiększa skuteczność agenta.
    \item[\textbf{H2}] Jednoznaczne, opisowe nazwy elementów interaktywnych i~pól formularzy zwiększają skuteczność realizacji zadań.
    \item[\textbf{H3}] Poprawne zastosowanie mechanizmów dostępności (etykiet, ról i~atrybutów ARIA) zwiększa skuteczność również w~przypadku agenta multimodalnego, jakim jest agent Antigravity.
    \item[\textbf{H4}] Uproszczenie niestandardowych komponentów i~zapewnienie jawnych komunikatów o~stanie zmniejsza liczbę nieudanych zadań oraz fałszywych deklaracji sukcesu.
    \item[\textbf{H5}] Wpływ modyfikacji interfejsu jest większy w~przypadku słabszych modeli językowych wykorzystywanych przez agenta.
\end{itemize}

Hipoteza H5 jest szczególnie istotna z~praktycznego punktu widzenia: jeżeli zostanie potwierdzona, oznaczać będzie, że odpowiednio zaprojektowana aplikacja pozwala na wykorzystanie tańszych i~szybszych modeli przy zachowaniu akceptowalnej skuteczności.

% =====================================================================
\subsection{Podsumowanie projektu}
\label{sec:podsumowanie-projektu}
Zaprojektowane środowisko łączy elementy, które w~istniejących rozwiązaniach występowały osobno: realistyczne aplikacje biurowe uruchamiane w~kontrolowanym środowisku (jak w~WebArena), dziedzinę zadań biurowych (jak w~WorkArena), powiązanie błędów agenta z~konstrukcją komponentów (jak w~WebSuite) oraz dbałość o~izolację zmiennych eksperymentu (jak w~BrowseComp-Plus). Elementem wyróżniającym jest uczynienie zmienną eksperymentu samej aplikacji, a~nie agenta, oraz wykorzystanie gotowego, komercyjnego agenta przeglądarkowego w~miejsce agenta budowanego przez badacza. Głównym produktem inżynierskim pracy jest środowisko wraz z~zestawem scenariuszy i~mechanizmem weryfikacji, umożliwiające rzetelny pomiar skuteczności agenta Google Antigravity niezależnie od tego, jak duże okażą się różnice pomiędzy wariantami interfejsu. Wyniki eksperymentów posłużą do sformułowania zaleceń projektowych dotyczących aplikacji przyjaznych agentom.

% TODO: kolejne podrozdziały (implementacja) zostaną dodane po realizacji środowiska.


\newpage
% Tymczasowo wyłączone (wersja częściowa: literatura + projekt); pliki zawierają jeszcze tekst wzorcowy formatki.
% \section{Opis przeprowadzonych testów, ich zakres i wyniki}
\label{chap:testy}
Przeprowadzenie kompleksowych testów jest nieodłącznym elementem każdej pracy naukowej i inżynierskiej. Dzięki testom możemy weryfikować skuteczność, efektywność oraz bezpieczeństwo opracowanego rozwiązania. Metody i rodzaje ewentualnych mierzalnych testów opisano w rozdziale ,,Cel i zakres pracy''. Rozdział ten zależeć będzie od podjętego tematu i zaleceń promotora. Poniższe podrozdziały są tylko przykładem.

\subsection{Zakres testów}

W ramach badań przeprowadzono trzy główne kategorie testów\ldots

\subsection{Wyniki testów}

Na początek, w krótkim wstępie, powinniśmy przedstawić charakterystykę testów oraz przypomnieć czytelnikowi cele i założenia, które zostały wcześniej opisane w rozdziale poświęconym metodologii testowania. Następnie warto przedstawić specyfikację środowiska testowego, uwzględniając zarówno specyfikację sprzętową, jak i programową, a także szczegółowe wersje używanego oprogramowania czy systemów operacyjnych.

Kolejnym krokiem jest podsumowanie wszystkich przeprowadzonych testów, wskazując, ile z nich zakończyło się sukcesem, a ile niepowodzeniem. Szczegółowe wyniki powinny zawierać informacje na temat funkcjonalności, wydajności oraz bezpieczeństwa. Jeśli przeprowadzono testy użyteczności, warto również dodać informacje na temat feedbacku od użytkowników oraz napotkanych trudności.

Niezbędne są też wizualizacje: wykresy, diagramy czy tabele, które pomogą czytelnikowi zrozumieć i zinterpretować uzyskane dane. Po prezentacji wyników ważna jest ich głęboka analiza, interpretacja w kontekście wcześniej postawionych celów oraz omówienie ewentualnych odchyleń od oczekiwań. Pamiętaj jednak, że sam wykres to za mało. Należy wskazać czytelnikowi, które aspekty uzyskanych wyników potwierdzają tezy pracy.

\section{Analiza aspektów prawnych}

W kontekście pracy inżynierskiej, analiza aspektów prawnych odgrywa kluczową rolę. Oprogramowanie, protokoły, algorytmy i różne metody, które zostały opracowane, podlegają różnym normom prawnym, które muszą być rozważone, aby zapewnić legalność oraz etyczność opracowanego rozwiązania.

\subsection{Prawo autorskie}

Wszystkie skrypty, kody źródłowe oraz dokumentacje stworzone w ramach tej pracy są chronione prawem autorskim. Jest to kluczowe, aby zabezpieczyć prawa własności intelektualnej autora oraz zapewnić odpowiednie licencjonowanie. Należy zastanowić się nad odpowiednią licencją, która umożliwi innym korzystanie z opracowanego rozwiązania, jednocześnie chroniąc prawa autora.

\subsection{Ochrona danych}

Jeśli w pracy zostały wykorzystane jakiekolwiek dane, szczególnie te dotyczące osób fizycznych, istotne jest przestrzeganie przepisów związanych z ochroną danych osobowych. W kontekście europejskim, należy zwrócić uwagę na zgodność z Rozporządzeniem o Ochronie Danych Osobowych (RODO).

\subsection{Prawo patentowe}

Jeśli opracowane rozwiązania są nowatorskie i mają potencjał komercyjny, warto rozważyć możliwość ubiegania się o patent. Chociaż proces patentowania w obszarze informatyki może być skomplikowany, chroni on przed nieautoryzowanym wykorzystaniem wynalazku.

\subsection{Licencjonowanie oprogramowania}

W zależności od charakteru rozwiązania oraz planowanej dystrybucji, istotne jest wybór odpowiedniej licencji. Może to być licencja open-source, która pozwala na szerokie wykorzystanie kodu przez społeczność, lub licencja komercyjna, która ogranicza dostęp do kodu i jego wykorzystanie.

\subsection{Etyka w nauce}

Oprócz ścisłych przepisów prawnych, ważne jest również przestrzeganie ogólnych zasad etyki w nauce. Obejmuje to uczciwość, transparentność oraz odpowiedzialność za stworzone rozwiązania i ich ewentualne wpływy na społeczeństwo.

Wnioski z niniejszego rozdziału mają na celu podkreślenie znaczenia aspektów prawnych w~procesie tworzenia nowych rozwiązań w dziedzinie informatyki stosowanej. Prawo w tej dziedzinie jest skomplikowane i dynamicznie się zmienia, dlatego istotne jest ciągłe monitorowanie aktualnych przepisów i dostosowywanie do nich opracowywanych rozwiązań.



\newpage
% Tymczasowo wyłączone (wersja częściowa: literatura + projekt); pliki zawierają jeszcze tekst wzorcowy formatki.
% 
\section*{Wnioski}
\addcontentsline{toc}{section}{Wnioski}
Na zakończenie rozdziału, w części poświęconej wnioskom, należy podsumować najważniejsze ustalenia i, jeśli to możliwe, zaproponować kierunki dalszej pracy, optymalizacji lub poprawek. Jeśli istnieją inne, podobne rozwiązania, warto także dodać porównanie uzyskanych wyników z tymi dostępnymi na rynku lub opisanymi w literaturze. Ostatecznie, najważniejsze jest, aby prezentacja wyników była przejrzysta, czytelna i konkretna.

\section*{Podsumowanie}
\addcontentsline{toc}{section}{Podsumowanie}
Rozdział ten stanowi kluczowy punkt pracy inżynierskiej, w którym wszystkie istotne elementy projektu i analizy zostają skumulowane i przedstawione w sposób zintegrowany. Jest to miejsce, gdzie podkreślana jest praktyczna wartość stworzonego rozwiązania technicznego oraz to, w~jaki sposób odpowiada ono na postawione wcześniej problemy inżynierskie.
W pierwszej kolejności, autor powinien podkreślić, czy cele pracy zostały osiągnięte i w jakim stopniu rozwiązanie techniczne spełnia założone wymagania.
Następnie ważne jest przedstawienie głównych osiągnięć projektowych, skupiając się na tym, co nowatorskiego wnosi prezentowane rozwiązanie i jakie korzyści może przynieść w praktyce.
Autor powinien również odnieść się do ewentualnych trudności napotkanych w trakcie realizacji projektu, omawiając wyzwania i~sposoby ich pokonywania.
W kontekście inżynierskim, warto porównać opracowane rozwiązanie z~dostępnymi na rynku alternatywami, wskazując jego mocne strony oraz potencjalne obszary do dalszej optymalizacji.
Praktyczne zastosowanie stworzonego rozwiązania powinno być podkreślone, omawiając potencjalne korzyści dla firm, instytucji lub użytkowników końcowych.
Zakończenie powinno wskazywać na potencjalne kierunki dalszych prac nad rozwiązaniem, co można jeszcze udoskonalić, jakie funkcje dodać czy w jakich obszarach zastosować wynikowe rozwiązanie inżynierskie.


\newpage

\section*{Wyjaśnienie skrótów używanych w pracy}
\addcontentsline{toc}{section}{Wyjaśnienie skrótów używanych w pracy}
\begin{description}[leftmargin=2.8cm, labelwidth=2.5cm, itemsep=0pt]
    \item[API] interfejs programistyczny aplikacji (ang. \textit{Application Programming Interface})
    \item[ARIA] zestaw ról, stanów i właściwości opisujących komponenty interfejsu dla technologii wspomagających (ang. \textit{Accessible Rich Internet Applications}, WAI-ARIA)
    \item[AWI] agentowy interfejs internetowy (ang. \textit{Agentic Web Interface})
    \item[CRM] system zarządzania relacjami z klientami (ang. \textit{Customer Relationship Management})
    \item[DOM] obiektowy model dokumentu (ang. \textit{Document Object Model})
    \item[HTML] hipertekstowy język znaczników (ang. \textit{HyperText Markup Language})
    \item[HTTP] protokół przesyłania dokumentów hipertekstowych (ang. \textit{Hypertext Transfer Protocol})
    \item[LLM] duży model językowy (ang. \textit{Large Language Model})
    \item[MCP] protokół kontekstu modelu (ang. \textit{Model Context Protocol})
    \item[TSR] wskaźnik skuteczności realizacji zadań (ang. \textit{Task Success Rate})
    \item[W3C] organizacja standaryzująca technologie internetowe (ang. \textit{World Wide Web Consortium})
    \item[WCAG] wytyczne dotyczące dostępności treści internetowych (ang. \textit{Web Content Accessibility Guidelines})
\end{description}

% Wykaz literatury
\newpage

\printbibliography

% Inne spisy

\newpage
\pagestyle{fancy}
\listoffigures
\clearpage
\newpage
\listoftables
\newpage
\listofmyequations
\newpage
\renewcommand{\lstlistlistingname}{Spis kodów}
\lstlistoflistings

\newpage
% Tymczasowo wyłączone (wersja częściowa: literatura + projekt); pliki zawierają jeszcze tekst wzorcowy formatki.
% \section*{Wykaz symboli i oznaczeń}
% Treść wykazu symboli i oznaczeń
\subsection*{Symbole łacińskie}

\nomenclature{$a$}{Przykładowy symbol}
\nomenclature{$x$}{Inna zmienna}

\subsection*{Symbole greckie}

\nomenclature{$\alpha$}{Pierwsza litera greckiego alfabetu}
\nomenclature{$\beta$}{Druga litera greckiego alfabetu}

\subsection*{Indeksy i wskaźniki}

\nomenclature{$i$}{Indeks dla zmiennej $x$}
\nomenclature{$\mathrm{max}$}{Wskaźnik oznaczający wartość maksymalną}

\subsection*{Skróty}

\nomenclature{CPU}{Central Processing Unit}
\nomenclature{RAM}{Random Access Memory}

\printnomenclature



\newpage
\section*{Wykaz używanych skrótów}
% Treść wykazu używanych skrótów
\newpage
\section*{Wykaz definicji}
% Treść wykazu definicji
\newpage
\section*{Wykaz załączników}
% Treść wykazu załączników

\end{document}
